% Benutzerhandbuch — rtheatflow (Echtzeit-Simulation thermischer Wärmenetze)
% Übersetzen mit:  latexmk -pdf Benutzerhandbuch.tex
% Aufbau und LaTeX-Konventionen gespiegelt vom Schwesterprojekt rtpowerflow (netzsim).
\documentclass[11pt,a4paper,parskip=half]{scrreprt}
\usepackage[T1]{fontenc}\usepackage[utf8]{inputenc}\usepackage{lmodern}
\usepackage[ngerman]{babel}\usepackage{microtype}\usepackage{amsmath}
\usepackage{booktabs}\usepackage{longtable}\usepackage{tabularx}\usepackage{xcolor}
\usepackage{enumitem}\usepackage{listings}\usepackage{graphicx}\graphicspath{{img/}}
\usepackage[colorlinks=true,linkcolor=blue!50!black,urlcolor=blue!50!black]{hyperref}
\definecolor{shell}{gray}{0.96}\definecolor{bsp}{RGB}{0,84,159}
\lstset{basicstyle=\ttfamily\small,backgroundcolor=\color{shell},frame=single,framerule=0pt,
  xleftmargin=6pt,xrightmargin=6pt,aboveskip=8pt,belowskip=8pt,breaklines=true,
  literate={ü}{{\"u}}1 {ä}{{\"a}}1 {ö}{{\"o}}1 {ß}{{\ss}}1 {Ü}{{\"U}}1 {Ä}{{\"A}}1 {Ö}{{\"O}}1
    {→}{{$\rightarrow$}}1 {Σ}{{$\Sigma$}}1 {≈}{{$\approx$}}1 {−}{{$-$}}1 {Δ}{{$\Delta$}}1 {·}{{$\cdot$}}1}
\newcommand{\ui}[1]{\textsf{\bfseries #1}}
\newcommand{\api}[1]{\texttt{#1}}
% Grad Celsius ohne Zusatzpaket (das gegebene Präambel-Set enthält kein textcomp):
\newcommand{\gC}{\,\ensuremath{{}^{\circ}\mathrm{C}}}
\emergencystretch=3em
\newenvironment{beispiel}[1][Beispiel]
  {\par\medskip{\color{bsp}\hrule height 0.8pt}\nopagebreak\smallskip
   \noindent\textbf{\textsf{\textcolor{bsp}{#1}}}\quad\ignorespaces}
  {\par\nopagebreak\smallskip{\color{bsp}\hrule height 0.8pt}\par\medskip}
\title{rtheatflow\\\Large Benutzerhandbuch}
\subtitle{Echtzeit-Simulation thermischer Wärmenetze\\für Lehre und Forschung}
\author{}
\date{Stand: Juli 2026}

\begin{document}

\maketitle

% Entstehungs- und Lizenzhinweis (Impressum)
\begin{center}
\fbox{\parbox{0.86\textwidth}{\small
\textbf{Hinweis zur Entstehung:} Der Quellcode von rtheatflow sowie dieses
Handbuch wurden mit einem \textbf{KI-Coding-Agenten} (Claude Code, Anthropic)
erstellt --- unter menschlicher Anleitung: Anforderungen, fachliche
Entscheidungen, Review und die Live-Verifikation gegen das laufende System
stammen von Menschen. Die Wärmenetz-Physik ist gegen den
\textbf{IBPSA-DESTEST-Cross-Tool-Benchmark} und gegen \textbf{reale
Verbier-Messdaten} validiert (Kapitel~\ref{ch:referenznetze}).
Quellcode und Handbuch stehen unter der MIT-Lizenz; die mitgelieferten
Referenzdatensätze behalten ihre eigenen Lizenzen (IBPSA DESTEST: modifizierte
BSD-3; pandapipes-Beispielnetze: BSD-3; OpenDHN/Verbier: CC~BY~4.0).
Projektseite: \url{https://github.com/markisbell/rtheatflow}}}
\end{center}

\tableofcontents

% ===================================================================
\chapter{Einführung}
\label{ch:einfuehrung}

\section{Was ist rtheatflow?}

rtheatflow ist eine Lehr- und Forschungsplattform für thermische
\emph{Wärmenetze} (Fernwärme). Sie löst fortlaufend die
Thermo-Hydraulik eines Zweileiter-Fernwärmenetzes auf Basis von
\href{https://github.com/e2nIEE/pandapipes}{pandapipes~0.14.0} und stellt die
Ergebnisse live im Browser dar --- auf einer echten OpenStreetMap-Karte, in der
sich Vorlauf- und Rücklauftemperaturen, Massenströme, Drücke und Wärmeverluste
Minute für Minute entwickeln. Anlagen (Erzeuger, Pufferspeicher, Einspeiser,
Bypässe, Abnehmer) lassen sich im laufenden Betrieb platzieren; eine Heizkurve
und eine Schlechtpunktregelung schließen die Betreiberschleife.

rtheatflow ist das \emph{Wärme-Schwesterprojekt} des Stromnetzsimulators
\href{https://github.com/markisbell/rtpowerflow}{rtpowerflow} (netzsim) und
spiegelt dessen Architektur --- Elektrizität wird konsequent auf Wärme
übersetzt: Spannung wird zu Temperatur, Strom zu Massenstrom, Wirkleistung zu
Wärmestrom, die Leitung zur Trasse, der Trafo zur Heizzentrale.

Das didaktische Kernkonzept ist die Trennung von \textbf{drei Sichten} auf
dasselbe Netz.

\section{Die drei Sichten}
\label{sec:dreisichten}

\begin{description}
  \item[\ui{Realität}] Der vollständige, physikalisch gelöste Zustand jeder
    Trasse, jedes Knotens, jedes Abnehmers --- das, was im Netz
    \emph{tatsächlich} passiert. In der Wirklichkeit kennt ihn niemand.
  \item[\ui{Gemessen}] Nur das, was platzierte Messgeräte
    (Wärmemengenzähler, T/p-Sensoren) tatsächlich erfassen --- und nur in dem
    Zeitraster, das der gewählte Messmodus liefert. Alles andere ist grau und
    unbekannt (\glqq unbekannt ist unbekannt\grqq). Das ist der Standardblick
    des Betreibers.
  \item[\ui{Schätzung}] Das, was ein Betreiber aus seinen wenigen Messwerten
    \emph{berechnen} kann: ein zweites Netzmodell, angetrieben allein von
    Messwerten und Erwartungsprofilen --- inklusive des Fehlers, den diese
    Rechnung gegenüber der Wahrheit macht.
\end{description}

Der intellektuelle Anspruch, den die Plattform im Code verteidigt: \emph{Die
Sichten Gemessen und Schätzung dürfen ausschließlich Informationen enthalten,
die ein realer Betreiber tatsächlich hätte.} Keine Ebene oberhalb der Realität
darf um die Projektion herum an die Wahrheit greifen --- und das ist kein
Vorsatz, den man sich merken muss, sondern durch Tests (die
\glqq Ehrlichkeits-Stolperdrähte\grqq, Abschnitt~\ref{sec:stolperdraehte}) und
zur Laufzeit durch den strikten Modus (Abschnitt~\ref{sec:strikt}) erzwungen.

\begin{beispiel}[Typische Lehrfragen, die sich damit beantworten lassen]
\begin{itemize}[nosep]
  \item Was passiert mit dem Schlechtpunkt-$\Delta p$, wenn morgens alle
    Übergabestationen gleichzeitig aufheizen?
  \item Wie ändert sich die Verlustquote, wenn man von einem 3G-Netz
    (110/60) auf ein 4G-Netz (70/40) umschaltet --- und was kostet das an
    Pumpenstrom?
  \item Wie viele Wärmemengenzähler braucht die Zustandsschätzung, um eine
    Rücklauf-Anomalie an einer einzelnen Station überhaupt zu \emph{sehen}?
  \item Kappt ein 100-kWh-Pufferspeicher die morgendliche Wärmespitze
    wirksam --- und wann hört er von selbst auf zu laden?
\end{itemize}
\end{beispiel}

\section{Die drei Anwendungen}

Die Plattform besteht aus drei Anwendungen, die gemeinsam über eine einzige
\texttt{docker-compose.yml} starten:

\begin{enumerate}
  \item \textbf{rtheatflow} --- der Simulationskern: eine FastAPI-Anwendung
    (\api{src/}), die ein Netz samt Tagesprofilen lädt und fortlaufend die
    Thermo-Hydraulik löst; jedes Ergebnis ist per REST abrufbar und wird
    zusätzlich über einen WebSocket an alle verbundenen Browser gepusht
    (Backend, \textbf{Port~8001}).
  \item \textbf{Web-Oberfläche} --- die React-Anwendung im Browser
    (\api{ui/}, deutschsprachig, oben rechts auf EN umschaltbar), über die alle
    Funktionen dieses Handbuchs bedient werden (Vite-Entwicklungsserver
    \textbf{Port~5174}, in Docker per nginx auf \textbf{Port~8081}).
  \item \textbf{Visualisierungs-Stack} --- ein Collector holt jeden gelösten
    Schritt per REST ab und schreibt ihn nach InfluxDB; ein vorkonfiguriertes
    Grafana-Dashboard zeigt die Langzeitverläufe (Kapitel~\ref{ch:grafana},
    Grafana \textbf{Port~3001}, InfluxDB \textbf{Port~8087}).
\end{enumerate}

\paragraph{Das Geschwister-Portschema.} Die Ports 8001/5174 (und 8081/8087/3001
im Docker-Stack) sind \emph{mit Absicht} gewählt: netzsim/rtpowerflow belegt
8000/5173 (bzw.\ 8080/8086/3000). So laufen beide Plattformen --- Strom und
Wärme --- parallel auf derselben Maschine, ohne sich die Ports streitig zu
machen.

\section{Das Zeitmodell}
\label{sec:zeitmodell}

Ein Simulationstag besteht aus \textbf{1440 Schritten zu je einer Minute}. Für
jeden Schritt werden Lasten, Wetter, Heizkurve und Anlagen auf die Werte des
Zeitschritts gesetzt und ein vollständiger, eingeschwungener
Thermo-Hydraulik-Zustand gelöst (auf üblichen Rechnern warm etwa 20--35~ms, siehe
Kapitel~\ref{ch:referenznetze}). Die Simulation läuft im \emph{beschleunigten
Takt}: pro \ui{Schrittdauer} realer Sekunden schreitet die Simulationsuhr eine
Minute voran.

\begin{center}
\begin{tabular}{@{}lll@{}}
\toprule
Schrittdauer & Beschleunigung & ein Simulationstag dauert \\
\midrule
1{,}0 s (Vorgabe) & 60$\times$ & 24 Minuten \\
0{,}5 s & 120$\times$ & 12 Minuten \\
0{,}1 s & 600$\times$ & 2{,}4 Minuten \\
\bottomrule
\end{tabular}
\end{center}

Nach Schritt~1439 (23:59~Uhr) springt die Simulation auf Schritt~0 zurück,
zählt den Tag hoch und wiederholt die Tagesprofile. Schrittindex und Uhrzeit
rechnen sich direkt ineinander um: $\text{Schritt} = h \cdot 60 + \mathit{min}$.

\begin{beispiel}
07:20~Uhr entspricht Schritt $7 \cdot 60 + 20 = 440$. Wer per API an die
morgendliche Aufheizspitze springen will, ruft daher
\api{POST /control/seek} mit \texttt{\{"step": 440\}} auf --- die Oberfläche
erledigt dasselbe über den Zeit-Regler in der Steuerleiste.
\end{beispiel}

\section{Validierte Physik}
\label{sec:validierung}

Ein Lehrwerkzeug ist nur so glaubwürdig wie seine Physik. Die Thermo-Hydraulik
von rtheatflow ist deshalb gegen \emph{externe} Wahrheit validiert, nicht gegen
frühere eigene Ausgaben: gegen den \textbf{IBPSA-DESTEST}-Cross-Tool-Benchmark
(mehrere unabhängige Gebäudesimulations-Werkzeuge rechnen dasselbe kleine Netz
und veröffentlichen ihre Ergebnisse --- die Streuung zwischen den Werkzeugen ist
selbst die Toleranz) und gegen \textbf{reale Messdaten} des vermaschten
Alpennetzes Verbier (150 Übergabestationen). Als Plausibilitätsprüfung dient
zusätzlich das reale Ortsnetz Schutterwald. Die vollständigen Zahlen, die
akzeptierten Modellunterschiede und die Reproduktionskommandos stehen in
Kapitel~\ref{ch:referenznetze}. Kurzfassung: Die Vorlauftemperaturen der
150~Verbier-Stationen werden im Median auf \textbf{0{,}83~K} genau getroffen,
ohne dass diese Temperaturen jemals Eingabe wären.

\section{Weiterführende Dokumentation}
\label{sec:weiterdoku}

Dieses Handbuch ist die bedienungsorientierte Einführung. Für die Tiefe gibt es
im Verzeichnis \api{docs/} fünf technische Deep-Dives (englisch, gegen den Code
mit \texttt{Datei:Zeile}-Ankern geschrieben):
\api{ARCHITECTURE.md} (System, Wireformat, Solverpolitik, Beobachtbarkeit,
quasistatische Ehrlichkeit), \api{PHYSICS\_AND\_CONTROLS.md} (pandapipes-Modell,
Retry-Ladder, Verluste, Heizkurve, Schlechtpunkt, Anlagen),
\api{OBSERVABILITY.md} (die drei Ebenen, der Vorwärts-Beobachter, die
Stolperdrähte), \api{REFERENCE\_NETWORKS.md} (der Fünf-Dateien-Vertrag, alle
Netze) und \api{BENCHMARKS.md} (Rechenzeiten + Validierungsprotokoll).
\api{SPEC.md} ist die bindende Spezifikation, \api{docs/API.md} die generierte
Endpunktreferenz. Das vorliegende Handbuch wird zusätzlich unter \api{GET /manual}
direkt aus dem Repository als HTML ausgeliefert.

% ===================================================================
\chapter{Installation und Start}
\label{ch:installation}

\section{Windows (Entwicklung, ein Klick)}

Der einfachste Weg auf einem Windows-Rechner:

\begin{lstlisting}
start_rtheatflow.bat
\end{lstlisting}

Der Starter erkennt bereits laufende Server (und startet sie nicht doppelt),
startet Backend (Port~8001) und Vite-UI (Port~5174) je in einer eigenen Konsole,
wartet auf \api{/health} und öffnet den Browser. Der \emph{allererste}
Rechenschritt kompiliert die numba-Kerne vor (Kaltstart, bis $\approx$60~s,
Abschnitt~\ref{sec:kaltstart}); solange zeigt der Starter \glqq Warte auf das
Backend\grqq. \textbf{Beenden:} \texttt{stop\_rtheatflow.bat} schließt beide
Serverfenster samt eventuell verwaister Hintergrundprozesse dieses Repos ---
ein parallel laufendes netzsim/rtpowerflow (Ports 8000/5173) bleibt unberührt.

Einmalige Vorbereitung (die virtuelle Umgebung anlegen):

\begin{lstlisting}
py -3 -m venv .venv
.venv\Scripts\pip install -r requirements.txt
\end{lstlisting}

\texttt{ui\textbackslash node\_modules} wird beim ersten Start automatisch per
\texttt{npm install} nachgezogen.

\section{Lokaler Start (Entwicklung, manuell)}

Wer die beiden Prozesse selbst in der Hand haben will, startet sie getrennt:

\begin{lstlisting}
:: Backend (Port 8001)
set PYTHONPATH=src
.venv\Scripts\python -m rtheatflow.main

:: UI (Port 5174; leitet /api und /ws an 127.0.0.1:8001 weiter)
cd ui && npm run dev
\end{lstlisting}

Der Vite-Server bindet bewusst \texttt{127.0.0.1} (nicht \texttt{localhost}) als
Proxy-Ziel und läuft mit \texttt{strictPort} --- er scheitert lieber laut, als
still auf einen anderen Port auszuweichen und dann ins falsche Backend zu
proxen (Abschnitt~\ref{sec:fehler} erklärt die Windows-IPv6-Falle dahinter).

\section{Docker Compose (Gesamtstack)}

Der komplette Stack --- Backend, UI, InfluxDB, Collector und Grafana --- startet
mit einem Befehl:

\begin{lstlisting}
docker compose up --build
\end{lstlisting}

{\small
\begin{center}
\begin{tabular}{@{}lll@{}}
\toprule
Dienst & Adresse & Inhalt \\
\midrule
ui       & \url{http://localhost:8081} & Karten-Oberfläche (nginx) \\
backend  & \url{http://localhost:8001} & REST + WebSocket; Swagger unter \api{/docs}, \\
         &                             & Monitor unter \api{/}, Handbuch unter \api{/manual} \\
grafana  & \url{http://localhost:3001} & Fernwärme-Dashboard (admin / admin) \\
influxdb & \url{http://localhost:8087} & Zeitreihenspeicher (Entwicklungs-Zugangsdaten) \\
\bottomrule
\end{tabular}
\end{center}}

Das Backend-Image bäckt das mitgelieferte \api{data/}-Verzeichnis ein und läuft
eigenständig; das Volume \api{./data} macht Aufzeichnungen und importierte Netze
dauerhaft. Alle Zugangsdaten im Compose-File sind \textbf{Entwicklungs-Vorgaben}
(Lehrwerkzeug, bewusst ohne Authentifizierung) --- vor einem Einsatz außerhalb
einer vertrauenswürdigen Maschine ändern.

\section{Kaltstart: der numba-JIT-Aufschlag}
\label{sec:kaltstart}

pandapipes verwendet numba mit \emph{bloßem} \texttt{@jit} (ohne
\texttt{cache=True}), also ohne Cache-Datei: der erste Lastfluss eines jeden
Prozesses kompiliert die Pipeflow-Kerne frisch --- etwa 7~s reine
JIT-Zeit, dazu kommt der Import von \texttt{pandapipes}/\texttt{numba}/\texttt{llvmlite}
selbst (4{,}9--12{,}1~s, stark schwankend). Das Backend versteckt das: unter
\texttt{AUTOSTART} läuft der erste Rechenschritt in einem Hintergrund-Thread,
sodass die JIT-Kompilierung das Binden des Sockets nicht blockiert. Praktisch
heißt das: Der \emph{erste} veröffentlichte Frame trägt eine \api{solve\_ms} in
den Tausenden --- das ist JIT, kein langsames Netz. Jeder Folgeschritt ist warm.

% ===================================================================
\chapter{Arbeitsablauf: Netz -- Lasten -- Live}
\label{ch:arbeitsablauf}

Die Oberfläche ist wie ein klassisches Desktop-Programm organisiert: Beim Start
öffnet sich die \emph{Live-Ansicht}, die Menüleiste am oberen Rand bündelt alle
Funktionen. Ganz oben in der Leiste:

\begin{description}[nosep]
  \item[\ui{Datei}] alles, was lädt, speichert oder exportiert: Szenarien
    speichern/laden/löschen (Kapitel~\ref{ch:szenarien}) sowie Aufzeichnung und
    Tage-Export (Kapitel~\ref{ch:export}). Läuft eine Aufzeichnung oder ein
    Export, erscheinen rechts in der Leiste klickbare Status-Chips
    (\ui{REC~n} / \ui{Export~p\,\%}).
  \item[\ui{Ansicht}] wählt die Farbebene der Karte (Vorlauf, Rücklauf,
    Geschwindigkeit, Differenzdruck --- Kapitel~\ref{ch:karte}).
  \item[\ui{Hilfe}] öffnet dieses Handbuch (\api{/manual}), die Swagger-Doku
    (\api{/docs}) und die Projektseite.
\end{description}

Rechts neben den Menüs sitzen zwei dauerhaft sichtbare Umschalter: der
\textbf{Bereich} \ui{Live~/ NetzStudio} (zwischen laufender Simulation und
Netz-Konfiguration wechseln) und der \textbf{Sicht-Umschalter}
\ui{Realität~/ Gemessen~/ Schätzung} (Kapitel~\ref{ch:beobachtbarkeit}) --- die
drei Sichten sind das Kernkonzept, deshalb ist die aktive Sicht immer ablesbar
und mit einem Klick wechselbar. Ganz rechts steht der Netz-Chip (Name~$\cdot$
Abnehmerzahl~$\cdot$ Trassenkilometer) und der DE/EN-Umschalter.

% TODO Screenshot: Menüleiste mit geöffnetem Datei-Menü, Sicht-Segment und Netz-Chip.

Der typische Arbeitsablauf ist ein Dreischritt --- Netz wählen, Lasten erzeugen,
live beobachten. Die ersten beiden Schritte leben im Bereich \ui{NetzStudio},
der in drei Spalten von links nach rechts durchlaufen wird.

\section{NetzStudio: die drei Spalten}
\label{sec:netzstudio}

\paragraph{Spalte 1 --- Netz wählen oder importieren.} Die linke Spalte listet
alle verfügbaren Netze, gruppiert in \ui{Bibliothek} (die mitgelieferten
Katalognetze) und \ui{Eigene Netze} (per \api{POST /networks/import}
importierte Fünf-Dateien-Bündel). Ein \ui{Netz importieren…}-Dateidialog nimmt
entweder die fünf Vertragsdateien zusammen oder ein einzelnes Bündel entgegen,
prüft sie und legt das Netz als \api{source="user"} ab. Der Katalog enthält
zwei Lehrnetze (\ui{Demo-Dorf}, \ui{Appendix~A}) und die
\textbf{Referenznetze} DESTEST (8/16/32~Gebäude), Schutterwald und Verbier ---
alle mit Herkunft und validierten Zahlen (Kapitel~\ref{ch:referenznetze}).

\paragraph{Spalte 2 --- Lastpolitik.} Die Mittelspalte weist den Abnehmerknoten
realistische Tagesprofile zu (die \emph{Lastpolitik}, SPEC~\S4.5). Grundlage ist
ein mitgelieferter Archetyp-Katalog: physikbasierte Jahresprofile aus
\textbf{demandlib} (VDI~4655, Raumheizung) und \textbf{OpenDHW} (stochastisches
Warmwasser). Drei Archetypen sind ausgeliefert:

{\small
\begin{center}
\begin{tabular}{@{}llr@{}}
\toprule
Archetyp & Beschreibung & Jahres-Raumwärme \\
\midrule
\texttt{EFH\_ALT\_4P}   & unsaniertes Einfamilienhaus, 4~Pers. & 18 MWh/a \\
\texttt{EFH\_SAN\_4P}   & KfW-saniertes EFH, 4~Pers.           & 9 MWh/a  \\
\texttt{MFH\_ALT\_10WE} & Mehrfamilienhaus, 10~Wohneinheiten   & 95 MWh/a \\
\bottomrule
\end{tabular}
\end{center}}

Einstellbar sind: der Gebäudemix (Archetyp-Auswahl), der Zuweisungsmodus
(\texttt{round\_robin} / \texttt{random}), ein \ui{Seed}, eine \ui{Skalierung},
das \ui{Tag-Perzentil} (0{,}9~=~kalter Wintertag, 0{,}5~=~Übergang,
0{,}05~=~Sommer), der \ui{Jitter} (deterministische $\pm$30-min-Verschiebung und
0{,}92--1{,}08-Amplitudenstreuung, damit gleiche Gebäude nicht synchron ziehen),
die DHW-Varianten und das Temperaturniveau \ui{3G}/\ui{4G}. Das
Temperaturniveau koppelt sowohl die Archetyp-Rücklauftemperaturen als auch die
Heizkurve (3G-Rückläufe sind wärmer als 4G, siehe
Tabelle~\ref{tab:treturn}).

\begin{table}[htbp]\centering
{\small
\begin{tabular}{@{}lrr@{}}
\toprule
Rücklauftemperatur & 3G-Preset & 4G-Preset \\
\midrule
\texttt{EFH\_ALT\_4P}   & 55\gC & 45\gC \\
\texttt{EFH\_SAN\_4P}   & 45\gC & 35\gC \\
\texttt{MFH\_ALT\_10WE} & 50\gC & 40\gC \\
\bottomrule
\end{tabular}}
\caption{Die vom Temperaturniveau gekoppelten Archetyp-Rücklauftemperaturen
(\texttt{TRETURN\_C} in \api{loadgen/}).}
\label{tab:treturn}
\end{table}

\paragraph{Spalte 3 --- Prüfen \& Anwenden.} Die rechte Spalte zeigt die
Vorschau: eine Zuweisungstabelle, Kennzahlen (Auslegungslast, Spitze,
Jahresenergie, Trassenlänge) und die \textbf{lineare Wärmedichte} [MWh/(m$\cdot$a)]
als eingefärbte Kachel gegen die Faustregel $\geq$1--1{,}5, dazu eine
Jahresdauerlinie als Sparkline gegen die Auslegungslast. \ui{Anwenden} lädt das
Netz in die Engine (der Live-Bereich öffnet sich neu, die Uhr läuft ab Tag~0 /
Schritt~0). Nichts an der Vorschau verändert die laufende Simulation --- erst
\ui{Anwenden} wirkt.

\begin{beispiel}[Warum die lineare Wärmedichte über die Wirtschaftlichkeit entscheidet]
Wählt man das \ui{Demo-Dorf}, meldet die KPI-Kachel eine lineare Wärmedichte
von etwa 0{,}39--0{,}46~MWh/(m$\cdot$a) --- \emph{unter} der Faustregel $\geq$1,
die Kachel färbt sich warnend. Das ist kein Fehler, sondern der Lehrpunkt: Ein
dünn besiedeltes Straßendorf mit langen Trassen und wenig Anschlussdichte ist
für ein klassisches Fernwärmenetz kaum wirtschaftlich --- die stehenden
Verluste der Trasse verteilen sich auf zu wenig verkaufte Wärme. Das reale
Ortsnetz Schutterwald erreicht mit seinem 50-m-Anschlussradius dagegen
1{,}34~MWh/(m$\cdot$a), genau im tragfähigen Band.
\end{beispiel}

\section{Die Live-Ansicht und die Steuerleiste}

Nach \ui{Anwenden} zeigt der Bereich \ui{Live} die Karte (links,
Kapitel~\ref{ch:karte}) und eine in der Breite ziehbare Seitenleiste (rechts)
mit einklappbaren Abschnitten: \ui{Übersicht}, \ui{Schlechtpunkt},
\ui{Messstellen}, \ui{Heizkurve}, \ui{Wetter} und den per Strg-Klick
angehefteten Element-Abschnitten. Am unteren Rand liegt die \textbf{Steuerleiste}:

\begin{description}[nosep]
  \item[Start/Pause] hält die Simulationsuhr an bzw.\ lässt sie weiterlaufen.
  \item[Zeit-Regler] springt an eine Tageszeit (Schritt $0\dots1439$).
  \item[Tag-Regler] erscheint nur bei mehrtägigen Wetterhorizonten
    (\texttt{n\_days>1}).
  \item[Schrittdauer] 0{,}1--1~s echte Zeit pro Simulationsminute.
\end{description}

Jeder Steuerbefehl synchronisiert die Anzeige aus der Antwort der Engine; die
Uhr selbst wird vom WebSocket getrieben und liegt auch bei 10~Schritten/s
höchstens einen Frame hinter dem Backend.

% ===================================================================
\chapter{Die Karte und die Farbebenen}
\label{ch:karte}

Die Karte ist eine rohe Leaflet-Karte (ohne react-leaflet) auf
CARTO/OSM-Kacheln. Sie ist \emph{einmal pro Topologie} aufgebaut --- eine
Polylinie je Trasse, ein Kreismarker je Abnehmer, Marker für Erzeuger, Anlagen
und Sensoren --- und wird pro Frame nur \emph{umgefärbt}
(\texttt{.setStyle()}), nie neu gebaut. So bleibt die Darstellung auch bei
hohem Tempo flüssig.

\section{Die vier Farbebenen}
\label{sec:ebenen}

Das Menü \ui{Ansicht} (oder der Umschalter auf der Karte) wählt, welche Größe
die Trassen und Marker einfärbt. Die Rampen sind \emph{domänenverankert} ---
nicht auf das Min/Max des Augenblicks normiert, sondern an physikalisch
bedeutsamen Schwellen fixiert, sodass eine Farbe immer dasselbe bedeutet:

\begin{description}
  \item[\ui{Vorlauftemperatur} (Standard)] warme Rampe (Gelb$\rightarrow$Rot);
    \glqq ganz heiß\grqq{} liegt \emph{exakt} bei der Auslegungs-Vorlauftemperatur
    der aktiven Heizkurve. Ein 4G-Netz mit 70\gC{} glüht bei 70 genau so wie
    ein 3G-Netz bei 110 --- die Rampe folgt dem Auslegungspunkt.
  \item[\ui{Rücklauftemperatur}] kühle Rampe (Blau$\rightarrow$Violett) über
    25--70\gC. Hohe Rückläufe (schlechte Auskühlung, das
    \glqq Low-$\Delta T$-Syndrom\grqq) springen violett heraus und fallen
    sofort auf.
  \item[\ui{Geschwindigkeit}] Grün (Ruhe) $\rightarrow$ Bernstein genau ab
    \textbf{1{,}5~m/s} $\rightarrow$ Rot ab \textbf{3~m/s} (Kapazitätsgrenze);
    das Vorzeichen wird ignoriert, weil Rücklaufrohre \glqq rückwärts\grqq{}
    fließen.
  \item[\ui{Differenzdruck}] Ampel an den Abnehmermarkern: Rot streng
    \emph{unterhalb} des Mindest-$\Delta p$ (\texttt{DP\_MIN\_BAR}, Vorgabe
    \textbf{0{,}5~bar}), Bernstein in einem 0{,}15-bar-Rand darüber, sonst Grün;
    die Trassen treten dabei zurück.
\end{description}

% TODO Screenshot: dieselbe Szene in allen vier Ebenen nebeneinander, mit Legende.

Jede Ebene bringt ihre eigene Farbleiste (Legende) mit, deren oberste Marke bei
der Vorlaufebene die live gelesene Auslegungstemperatur zeigt.

\section{Unbekanntes ist unbekannt}
\label{sec:unbekannt}

Das Leitprinzip der Farbgebung: \emph{Was nicht gemessen ist, wird als
unbekannt dargestellt} --- niemals in einer \glqq gesunden\grqq{} Farbe. Ein
Element ohne Messwert bekommt ein eigenes, gedämpftes Grau
(\texttt{UNOBSERVED} \texttt{\#39424f} für Marker, \texttt{UNOBSERVED\_LINE}
\texttt{\#2b323c} für Trassen) plus ein \texttt{6\,6}-Strichmuster, das keine
gesunde Rampe je erzeugen kann. In der Sicht \ui{Gemessen} tragen die Trassen
gar keine Zähler --- sie sind deshalb grundsätzlich grau-gestrichelt; die
Abnehmer werden nur aus \api{measurements.consumers} gefärbt. Ein
Standard-Zähler, der noch in seinem ersten 15-Minuten-Fenster steckt, erscheint
mit gestricheltem \texttt{3\,3}-Ring und gedämpfter Füllung (ehrlicher
Kaltstart, Kapitel~\ref{ch:beobachtbarkeit}).

\section{Popups und Element-Kontextmenü}

Ein einfacher Klick auf eine Trasse, einen Abnehmer oder Erzeuger öffnet ein
Popup mit den Live-Werten. Popups sind \emph{lebendig}: Der Inhalt liest bei
jedem Frame den neuesten Zustand, ein offenes Popup aktualisiert sich also
laufend --- in der Realitätssicht tickt \glqq Vorlauf~$\rightarrow$~Rücklauf\grqq,
in der Sicht Gemessen steht dort \glqq keine Messung --- unbekannt\grqq, ganz
ohne Neubinden.

\ui{Strg-Klick} heftet ein Element als eigenen Abschnitt in die Seitenleiste
(mit Live-Werten und Konfigurationsreglern; ein Pufferspeicher zeigt dort seinen
SoC-Balken). \ui{Rechtsklick} öffnet das Element-Kontextmenü --- die
Grammatik zum Platzieren und Konfigurieren von Anlagen und Messgeräten
(Kapitel~\ref{ch:anlagen} und~\ref{ch:beobachtbarkeit}).

% ===================================================================
\chapter{Beobachtbarkeit: Messen und Schätzen}
\label{ch:beobachtbarkeit}

Dies ist das Herzstück der Plattform. Ein realer Fernwärmebetreiber sieht nie
den wahren Zustand seines Netzes. Er sieht eine Handvoll Wärmemengenzähler und
Druck-/Temperatursensoren sowie die Instrumentierung der Heizzentrale --- und
aus diesem spärlichen Bild muss er entscheiden, wie hart die Pumpe läuft und wie
heiß der Vorlauf ist. In der Lücke zwischen wahrem und gemessenem Zustand leben
die Betriebsfehler: ein ausgehungerter Abnehmer am Netzende (der
\emph{Schlechtpunkt}), den kein Sensor abdeckt; ein hoher Rücklauf, der still
Pumpenstrom verschwendet; eine Störung, die keinen Alarm auslöst, weil kein
Instrument sie erfasst. rtheatflow macht diese Lücke sichtbar und
\emph{erzwingt} sie.

\section{Messgeräte platzieren}
\label{sec:messgeraete}

Drei Gerätearten stehen zur Verfügung. Zwei sind platzierbar (Rechtsklick auf
das Element $\rightarrow$ Kontextmenü, verwaltet im Seitenpanel-Abschnitt
\ui{Messstellen}); die dritte ist immer da:

\begin{description}
  \item[\ui{Wärmemengenzähler}] an einer Übergabestation (Abnehmer): liefert
    Leistung, Massenstrom, Vor- und Rücklauftemperatur sowie $\Delta p$
    (Kanäle \texttt{q\_kw}, \texttt{mdot\_kg\_per\_s}, \texttt{t\_supply\_c},
    \texttt{t\_return\_c}, \texttt{dp\_bar}).
  \item[\ui{T/p-Sensor}] an einem Trassenknoten: Druck und Temperatur in Vor-
    und Rücklauf des Knoten-Paares.
  \item[\ui{Erzeuger-SCADA}] an der Heizzentrale, \textbf{immer} vorhanden ---
    echte Heizwerke messen sich selbst. Sie ist Telemetrie, keine Messstelle,
    erscheint nie als Platzierung und degradiert nie in das 15-Minuten-Raster.
\end{description}

Trassen tragen bewusst \emph{keine} Messung --- Leitungsströme und Rohrverluste
bleiben unbekannt, solange die Wahrheit ausgeblendet ist. Gerade das
höchstbelastete Rohr ist also unsichtbar, wenn dort niemand misst.

\begin{beispiel}[Die Summe lügt nicht --- sie ist nur unvollständig]
Die Sicht \ui{Gemessen} zeigt in der Übersicht die \glqq gemessene
Nachfrage~($\Sigma$)\grqq: die Summe nur der Zähler, die \emph{gerade einen
Wert liefern}. Ein unbemessener Abnehmer trägt nichts bei, ein
Standard-Zähler im Kaltstart-Fenster trägt nichts bei. Diese Zahl ist also
ehrlich \emph{kleiner} als die wahre Nachfrage, sobald die Abdeckung
lückenhaft ist. Die SCADA-Kennzahlen der Zentrale (Einspeisung, Vorlauf,
Rücklauf, Pumpenstrom) sind dagegen belastbar, weil durch die Zentrale
\emph{alles} fließt. Diese Asymmetrie --- Summenmessung an der Zentrale
belastbar, Ortsauflösung im Netz null --- ist der Kern des
Beobachtbarkeitsproblems.
\end{beispiel}

\subsection*{Vorlagen (Presets)}

Der Abschnitt \ui{Messstellen} bietet vier Vorlagen, die die Platzierung
jeweils \emph{komplett ersetzen} (Sammelaktionen, keine additiven Änderungen):

\begin{description}[nosep]
  \item[\ui{Alle Abnehmer}] ein Zähler an jeder Übergabestation --- volle
    Abdeckung, die Vorgabe.
  \item[\ui{Nur Erzeuger}] nur das T/p-Paar an der Zentrale (SCADA ist ohnehin
    immer an).
  \item[\ui{Schlüsselstellen}] Zentrale + T/p an den Netzenden + ein Zähler am
    \emph{aktuell bekannten} Schlechtpunkt (aus dem letzten konvergierten
    Frame; vor dem ersten Rechenschritt ist keiner bekannt und keiner wird
    gesetzt --- ehrlich).
  \item[\ui{Alles entfernen}] keine Geräte --- Blindflug auf reiner
    Erzeuger-SCADA.
\end{description}

\subsection*{Zähler-Modus: voll oder 15 Minuten}
\label{sec:fidelity}

Ein einziger Modus-Schalter (\api{POST /measurements/mode}) gilt für alle
platzierten Geräte:

\begin{description}[nosep]
  \item[\ui{Voll}] jeder Kanal, jeder Simulationsschritt.
  \item[\ui{Standard (15 min)}] 15-Minuten-Fenstermittelwerte, ausgerichtet an
    der \emph{Simulationszeit} (Fenster $w = \text{Tick} // \text{Fensterschritte}$,
    tagesübergreifend stabil) --- das Verhalten eines deutschen
    Lastgangzählers. Der Kanal liest den Mittelwert des zuletzt
    \emph{abgeschlossenen} Fensters und liefert bis zur ersten Fenstergrenze
    nach dem Platzieren ehrlich \textbf{gar nichts} (\texttt{null}).
\end{description}

Der \textbf{ehrliche Kaltstart} ist die Pointe: Ein frisch gesetzter
Standard-Zähler zeigt \texttt{null} --- nicht null~kW, nicht den Momentanwert
--- bis sein erstes 15-Minuten-Fenster schließt. Ein Umschalten des Modus setzt
den Fensterzustand zurück (nie ein gefälschter Momentanwert aus einem alten
Akkumulator), ein entferntes Gerät verliert seinen Fensterzustand, ein
neu gesetztes startet wieder kalt.

\begin{beispiel}[Warum das 15-Minuten-Raster Spitzen verschluckt]
Zieht ein Abnehmer zwischen 07:02 und 07:12~Uhr zehn Minuten lang eine
Warmwasser-Spitze, zeigt ein \ui{Voll}-Zähler diese Spitze minutengenau. Ein
\ui{Standard}-Zähler liefert für das Fenster 07:00--07:15 nur den Mittelwert
--- rund $\tfrac{10}{15}$ der Spitze, ein Drittel ist \emph{in den Messdaten}
verschwunden (im Rohr floss es natürlich). Der Modus-Schalter macht den
Unterschied unmittelbar sichtbar --- und mit ihm die Frage, ob die
Schlechtpunktregelung dieser Spitze überhaupt folgen kann
(Abschnitt~\ref{sec:blindflug}).
\end{beispiel}

\section{Strikter Modus}
\label{sec:strikt}

Ist der Server im strikten Modus konfiguriert
(\api{RTHEATFLOW\_EXPOSE\_GROUND\_TRUTH=false}), verlässt die Wahrheit den
Server gar nicht erst. Der Strip passiert an genau \emph{einer} Stelle
(\texttt{StateStore.\_project}): Die vier Wahrheits-Schlüssel
(\texttt{junctions}, \texttt{pipes}, \texttt{consumers}, \texttt{summary})
werden entfernt und das Freitext-Feld \texttt{error} geleert. \api{/state},
\api{/history}, jeder WebSocket-Frame und (seit der Aufzeichnung) auch die
CSV-Pakete gehen durch diesen einen Pfad --- der strikte Modus lässt sich also
nicht durch einen vergessenen Endpunkt umgehen. Der Sicht-Umschalter fällt in
der Oberfläche automatisch von \ui{Realität} auf \ui{Gemessen} zurück und zeigt
den Hinweis \glqq Wahrheit ausgeblendet\grqq. Was bleibt: \texttt{measurements},
\texttt{observed\_summary}, die Betreiber-Konfiguration
(\texttt{producers}/\texttt{storages}/\texttt{controls}) --- und die
\texttt{estimated}-Ebene, denn sie ist aus Messwerten rekonstruiert und
gehört legitim dem Betreiber (Abschnitt~\ref{sec:schaetzung}). Damit lassen sich
Übungen abnahmesicher stellen (\glqq Schätzen Sie den Rücklauf am Netzende\grqq),
ohne dass die Lösung per API abgreifbar wäre.

\section{Der Vorwärts-Beobachter (die Schätzung)}
\label{sec:schaetzung}

pandapipes hat \emph{keinen} Zustandsschätzer --- nichts wie die
gewichtete Kleinste-Quadrate-Schätzung (WLS) von pandapower. Die Sicht
\ui{Schätzung} ist deshalb keine WLS-Rekonstruktion, sondern ein
\textbf{Vorwärts-Beobachter} (ein \glqq digitaler Zwilling\grqq): ein zweites,
tief kopiertes pandapipes-Netz, das mit derselben Physik (und derselben
Retry-Ladder) gelöst wird, aber \emph{nur mit Betreiberwissen} gefüttert:

\begin{itemize}[nosep]
  \item die SCADA-Vorlauftemperatur der Zentrale und der vom Betreiber gesetzte
    Pumpen-Sollwert (\texttt{plift});
  \item die Messwerte der platzierten Zähler, \emph{modus-treu} (ein
    Standard-Zähler im Kaltstart liefert \texttt{null} $\rightarrow$ der Prior
    übernimmt diesen Kanal);
  \item für alle \textbf{unbemessenen} Abnehmer ein \textbf{Erwartungsprofil}
    (Prior): das deterministische Archetyp-Profil, wettergestützt am Tag
    gewählt (eine Gradstunden-Lastprognose: erwartete Tagesenergie
    $\rightarrow$ nächster Archetyp-Tag), mit dem Mittel des Warmwassers über
    die stochastischen Varianten;
  \item die Anlagen-Sollwerte des Betreibers (Erzeuger, Speicher --- das ist
    Konfiguration, immer sichtbar);
  \item die wahre Außentemperatur inklusive Override (die Zentrale hat eine
    Wetterstation, und der Override ist der eigene Regler des Betreibers).
\end{itemize}

Entscheidend: Die Prioren lesen den \emph{unveränderlichen Planungsvertrag}
(\texttt{sim.inputs}), den Archetyp-Cache und die Platzierungs-Rezepte ---
\textbf{niemals} die veränderlichen Laufzeit-Arrays (\texttt{sim.profiles}), die
die stochastische Minutenwahrheit tragen. Eine zur Laufzeit injizierte Anomalie
kann deshalb prinzipiell nicht in den Prior lecken. Genau diese Disziplin
prüfen die Stolperdrähte (Abschnitt~\ref{sec:stolperdraehte}). Der Prior-Modus
ist über \api{POST /estimation/config} wählbar: \texttt{archetype} (Vorgabe, das
treuere Erwartungsprofil) oder das bewusst grobe \texttt{design}
(Auslegungslast~$\times$~Gradstundenfaktor, ohne Warmwasser).

\subsection*{Schätzgüte: die Innovation}
\label{sec:schaetzguete}

Die Übersicht zeigt in der Schätzsicht ein \textbf{Güte-Abzeichen}. Der
Fehler \texttt{estimated.error} ist die \emph{Innovation} des Beobachters:
$|\text{Zwilling} - \text{Messung}|$ an jeder \emph{bemessenen} Stelle ---
Rücklauftemperatur, Massenstrom und $\Delta p$ an Zählern; Rücklauf und
Paar-$\Delta p$ an T/p-Sensoren; Rücklauf und Massenstrom an der SCADA. Er wird
also gegen \emph{Messwerte} berechnet, nicht gegen die Wahrheit --- deshalb ist
er auch im strikten Modus gültig und sichtbar. Aggregiert erscheinen
$\max/\text{mean}$ von $|\Delta T_\text{Rücklauf}|$, $|\Delta\dot m|$ und
$|\Delta\Delta p|$, dazu das Alter der Schätzung (in Simulationsminuten, mit
Sequenznummer \texttt{\#seq}) und die Rechenzeit des Beobachters.

\begin{beispiel}[Volle Abdeckung kauft Gewissheit]
Mit der Vorlage \ui{Alle Abnehmer} im Modus \ui{Voll} ist der Zwilling an
jedem Abnehmer durch die gemessenen Randbedingungen (Leistung + Rücklauf) und
an der Zentrale durch die gemessene Vorlauftemperatur festgenagelt --- die
gesamte gekoppelte Lösung ist bestimmt. Die Schätzung trifft die Wahrheit dann
praktisch exakt (Einspeisung auf 0{,}1~kW, Massenstrom auf $10^{-3}$~kg/s), die
Innovation geht gegen null. Entfernt man Zähler, treten Prioren an die Stelle
von Sensoren, und die Innovation wächst --- diese monotone Beziehung ist das
quantitative Gesicht von \glqq Abdeckung kauft Gewissheit\grqq.
\end{beispiel}

\subsection*{Was eine Schätzung kostet (Drosselung)}

Der Beobachter ist ein zweites volles Netz --- wenn er rechnet, verdoppelt er
grob die Physikkosten des Schritts. Er rechnet deshalb \emph{nicht} jeden Tick.
Zwei Tore werden UND-verknüpft: das \textbf{Messraster} (im Standard-Modus mit
platzierten Geräten nur an den 15-Minuten-Fenstergrenzen --- dazwischen gibt es
keine neue Information) und eine \textbf{Wanduhr-Selbstdrosselung} (eine neue
Schätzung erst, nachdem \texttt{throttle\_factor}~$\times$~ihre eigene letzte
Rechenzeit verstrichen ist, Vorgabe 2{,}0). Zwischen den Aktualisierungen reitet
die letzte Schätzung auf jedem Frame mit, ehrlich mit \texttt{step}/\texttt{day}/\texttt{seq}
gestempelt, sodass die Oberfläche ihr Alter zeigen kann.

\begin{beispiel}[Sichtbare Rechenzeit auf dem vermaschten Netz]
Auf einem kleinen Lehrnetz kostet der Zwilling nur wenige zehn Millisekunden.
Auf dem vermaschten Verbier-Netz (1352~Knoten, Tier-2-Solver, $\approx$5~s je
Lösung) hält die Selbstdrosselung ihn davon ab, überhaupt jeden Tick zu
laufen: Die Schätzsicht zieht in größeren Sprüngen nach. Das ist kein Fehler,
sondern sichtbare Rechenzeit --- Beobachtbarkeit hat neben den Gerätekosten
auch Rechenkosten, und beide wachsen mit der Netzgröße.
\end{beispiel}

\section{Ehrlichkeits-Stolperdrähte}
\label{sec:stolperdraehte}

Die Stolperdrähte sind die Tests, die den Beobachtbarkeitsvertrag
\emph{falsifizierbar} machen. Jeder injiziert eine Lage, in der ein naiver
Schätzer schummeln würde, und prüft, dass er es nicht tut.

\begin{beispiel}[Eine unbemessene Anomalie bleibt unsichtbar]
Man bemisst nur Abnehmer~B und injiziert an Abnehmer~A (unbemessen) eine
Störung direkt in die Laufzeitwahrheit: verdoppelte Nachfrage und $+15$~K
Rücklauf. Die \ui{Realität} zeigt A prompt mit über 155~kW und über 68\gC{}
Rücklauf --- die Anomalie \emph{ist} in der Wahrheit. Die \ui{Schätzung} bleibt
beim Prior (rund 80~kW Plan, unter 57\gC), denn kein Sensor deckt A ab und die
Prioren lesen den unveränderlichen Vertrag, nicht die Laufzeit-Arrays.
Verschiebt man denselben Zähler auf A, \emph{propagiert} dieselbe Anomalie
sofort (die Schätzung zeigt über 155~kW). Zusammen beweisen beide Fälle: Die
Schätzung trägt exakt die Information, die die Sensorik liefert --- nicht mehr
und nicht weniger. \emph{So viele Wärmemengenzähler braucht die Schätzung, um
eine Rücklauf-Anomalie zu sehen: mindestens einen --- an genau dieser Station.}
\end{beispiel}

\begin{beispiel}[Ohne Zähler ist die Schätzung der Prior]
Unter der Vorlage \ui{Alles entfernen} (nur SCADA) ist jeder geschätzte
Abnehmerwert exakt gleich seinem Prior --- der Zwilling wird von nichts anderem
getrieben. Die Schätzgüte (z.\,B.\ die Innovation des Erzeuger-Massenstroms)
zeigt dann ehrlich, wie weit die Erwartung von der Wahrheit abweicht: Am
Demo-Dorf um 08:00~Uhr steigt die Massenstrom-Innovation von 0{,}000
(\ui{Alle Abnehmer}) über halbe Abdeckung bis 0{,}189~kg/s (\ui{Alles
entfernen}) --- die Degradation ist in genau der Zahl sichtbar, die die
Oberfläche anzeigt.
\end{beispiel}

\section{Der Blindflug der Schlechtpunktregelung}
\label{sec:blindflug}

Hier treffen sich Beobachtbarkeit und Regelung. Die Schlechtpunktregelung
(Kapitel~\ref{ch:physik}) liest \textbf{ausschließlich}
\texttt{observed\_summary.dp\_worst\_bar} --- das Minimum des $\Delta p$ über
die \emph{bemessenen} Abnehmer --- niemals die Wahrheit. Daraus folgt eine
Degradationsleiter:

\begin{enumerate}[nosep]
  \item \textbf{Kein verwertbarer $\Delta p$-Messwert} (Vorlagen \ui{Nur
    Erzeuger}/\ui{Alles entfernen}, oder alle Zähler noch im Standard-Kaltstart)
    $\rightarrow$ der Wert ist \texttt{null} $\rightarrow$ die Regelung
    \textbf{hält} die Förderhöhe. Halten \emph{ist} \glqq auf den Zentralen-$\Delta p$
    regeln\grqq: eine Konstant-$\Delta p$-Pumpe ohne Fernfühler.
  \item \textbf{Bemessen, aber der wahre Schlechtpunkt trägt keinen Zähler}
    $\rightarrow$ die Regelung regelt den besten \emph{gemessenen} Punkt (was
    ein realer Schlechtpunktregler auch tut) und der Frame setzt
    \texttt{controls.dp\_control.blind\_spot~=~true}.
\end{enumerate}

Das \texttt{blind\_spot}-Flag ist bewusste Meta-Information über die
\emph{Eignung der Sensorplatzierung}: Es benennt keinen Physikwert, nur einen
Booleschen Zustand darüber, ob die Sensorik den wahren Schlechtpunkt abdeckt ---
und bleibt deshalb auch im strikten Modus sichtbar (es teacht den Betreiber,
einen Zähler zu \emph{versetzen}, nicht der Zahl zu misstrauen). Der Abschnitt
\ui{Schlechtpunkt} zeigt zwei verschiedene Hinweistexte für die beiden Fälle.

\begin{beispiel}[Der wandernde Schlechtpunkt]
Mit der Vorlage \ui{Schlüsselstellen} liegt der Zähler dort, wo der
Schlechtpunkt beim letzten Rechenschritt \emph{bekannt} war --- nicht
notwendig dort, wo er \emph{jetzt} ist. Am Demo-Dorf pendelt der wahre
Schlechtpunkt unter der Warmwasser-Streuung zwischen zwei Abnehmern (C9~$\leftrightarrow$~C5).
Das \texttt{blind\_spot}-Flag geht genau dann auf \texttt{true}, wenn der
unbemessene C9 kritisch wird --- die Gemessen-Ebene meldet dann ehrlich den
bemessenen C5 bei einem \emph{höheren} $\Delta p$. Man sieht die Flagge kommen
und gehen: die Lektion vom wandernden Schlechtpunkt, live auf dem Bildschirm.
Das ganze Drei-Sichten-Modell existiert, um genau das sichtbar zu machen.
\end{beispiel}

% ===================================================================
\chapter{Physik und Regelungen}
\label{ch:physik}

Dieses Kapitel erklärt, was der Solver rechnet und was jeder Regler-Knopf tut.
Alle Formeln und Zahlen sind das, was der Code an \texttt{master} und seine
bekannte Antwort-Referenz tatsächlich produzieren.

\section{Das Vor-/Rücklaufmodell}

pandapipes ist ein stationärer Rohrströmungslöser. rtheatflow modelliert ein
Zweileiternetz --- heißer Vorlauf, kühler Rücklauf ---, indem der Netzbauer
\emph{jeden} Trassenknoten und \emph{jede} Trasse in ein \textbf{Paar} von
pandapipes-Elementen entfaltet, und zwar \emph{einmal} pro Konfiguration:

\begin{itemize}[nosep]
  \item ein Trassenknoten $\rightarrow$ zwei \texttt{junction}s
    (\texttt{<name>\_s} Vorlauf, \texttt{<name>\_r} Rücklauf), beide
    initialisiert auf die Vorlauftemperatur der Zentrale;
  \item eine Trasse $\rightarrow$ zwei \texttt{pipe}s (Vorlauf
    \texttt{from}$\rightarrow$\texttt{to}, Rücklauf
    \texttt{to}$\rightarrow$\texttt{from}, zur Zentrale hin umgekehrt);
  \item ein Abnehmer $\rightarrow$ ein \texttt{heat\_consumer}
    (Vorlauf-$\rightarrow$Rücklauf-Knoten), zieht \texttt{qext\_w};
  \item die Zentrale $\rightarrow$ genau eine \texttt{circ\_pump\_const\_pressure}
    (der Druck-Slack, Rücklauf$\rightarrow$Vorlauf).
\end{itemize}

Jeder Tick überschreibt nur die Elementwerte aus dichten
\texttt{[Elemente, Ticks]}-Arrays und löst neu --- das Netz wird nie neu gebaut.

\paragraph{Zwei Bau-Invarianten.} Erstens: Temperaturen sind \emph{Kelvin im
Solver, Grad Celsius auf dem Draht} (Rundung auf 6~Stellen, NaN/$\pm$Inf
$\rightarrow$ \texttt{null}). Zweitens: \texttt{text\_k} (die Erdreich-Temperatur
am Rohr) wird auf \emph{jedem} Rohr explizit gesetzt --- der Signatur-Standard
$0$ hieße 0~K Umgebung und würde die Verluste etwa vervierfachen. Jeder Tick
frischt \texttt{text\_k} aus der Erdreich-Zeitreihe auf.

\section{Die bidirektionale Lösung und die Retry-Ladder}

Temperaturgeregelte Abnehmer machen den Massenstrom von der \emph{ankommenden}
Temperatur abhängig (ein Abnehmer mit Rücklaufvorgabe muss bei kühlerem Vorlauf
mehr Massenstrom ziehen, um dieselbe Leistung zu liefern). Das koppelt
Hydraulik und Wärme, deshalb ist \texttt{mode="bidirectional"} die Vorgabe.
Konvergiert ein Schritt nicht sofort, geht der Löser eine vierstufige
\textbf{Retry-Ladder} von oben nach unten durch:

{\small
\begin{center}
\begin{tabular}{@{}clll@{}}
\toprule
Stufe & Aufruf & Zweck & Status \\
\midrule
1 & \texttt{bidirectional, iter=N}            & Normallösung        & \texttt{ok} \\
2 & \texttt{bidirectional, iter=N, alpha=0.5} & Halbschritt-Dämpfung & \texttt{ok} \\
3 & \texttt{bidirectional, iter=2N, alpha=0.2}& starke Dämpfung      & \texttt{ok} \\
4 & \texttt{sequential, iter=N}               & entkoppelt (Sollwerte nicht gehalten) & \texttt{degraded} \\
\bottomrule
\end{tabular}
\end{center}}

$N$ ist \api{RTHEATFLOW\_SOLVER\_ITER} (Vorgabe 100). \texttt{nonlinear\_method="automatic"}
wird bewusst \emph{nie} verwendet (es löst unter \texttt{bidirectional} in
0.14.0 einen Broadcast-Fehler aus). \textbf{Nichtkonvergenz ist Daten, kein
Absturz:} Scheitern alle Stufen, wird der letzte konvergierte Zustand
weiterverwendet und ein Frame mit \texttt{converged=false} /
\texttt{solver\_status="failed"} veröffentlicht --- die asyncio-Schleife lebt
weiter, ein Solverproblem erscheint als Feld und UI-Abzeichen, nie als
HTTP~500. Nach jedem konvergierten Schritt wird der gelöste Knotenzustand zur
nächsten Initialisierung (\emph{Warmstart}) --- das hält warme Lösungen bei
$\approx$20--35~ms.

\section{Abgeleitete Größen: Verluste und Einspeisung}

pandapipes liefert rohe Ergebnistabellen; die Plattform rechnet jede berichtete
Kennzahl selbst (in \texttt{collect\_physics}, geteilt von Wahrheit \emph{und}
Zwilling --- ein Formelsatz, nie zwei).

\begin{beispiel}[Warum die Verlustformel richtungsabhängig ist]
Der Verlust eines Rohres ist $q_\text{loss} = |\dot m| \cdot \bar c_p \cdot
(T_\text{ein} - T_\text{aus})$. Der Kniff: \emph{welches Ende ist der
Einlass?} Auf einem vermaschten Netz (Verbiers sechs Schleifen) kann ein Rohr
in beide Richtungen fließen, und die Rücklaufrohre fließen ohnehin
\glqq rückwärts\grqq{} zu ihrer \texttt{from}$\rightarrow$\texttt{to}-Definition.
Die Formel liest den Einlass aus dem Vorzeichen des Massenstroms
($T_\text{ein}=t_\text{from}$ falls $\dot m\geq0$, sonst $t_\text{to}$). Eine
naive Formel, die immer \texttt{t\_from} als Einlass nimmt, bekommt auf jedem
rückwärts fließenden Ast das Vorzeichen falsch und \textbf{überzählt den
Gesamtverlust um $\approx$20~\%} auf vermaschten Netzen.
\end{beispiel}

\begin{beispiel}[Einspeisung: $\dot m \cdot \bar c_p \cdot \Delta T$, nicht die Enthalpiespalte]
Die Einspeisung der Zentrale wird aus dem Ergebnis der Slack-Pumpe als
$\dot m \cdot \bar c_p \cdot \Delta T$ berechnet --- \emph{nicht} aus der Spalte
\texttt{res\_circ\_pump\_pressure.qext\_w}. Jene ist eine Enthalpieform, die
etwa $+4{,}7~\%$ zu hoch liest und die Energiebilanz nicht schließt: Auf der
bekannten Antwort-Referenz sind es \textbf{187{,}182~kW} ($\dot m\bar c_p\Delta T$)
gegen 195{,}965~kW (Rohspalte). Die Plattform berichtet den bilanzschließenden
Wert; \texttt{summary.balance\_err\_kw} ist der Ehrlichkeits-Stolperdraht pro
Schritt (auf der Referenz $-0{,}04~\%$).
\end{beispiel}

\section{Die Heizkurve}
\label{sec:heizkurve}

Die Heizkurve ist die Vorlauf-Temperaturregelung der Zentrale: kälter draußen
$\rightarrow$ heißerer Vorlauf (gleitender Betrieb). Sie folgt einer gleitenden
Kurve mit Formexponent $n$:
\begin{align*}
t_\text{flow}(t_\text{amb}) &= \text{clamp}\!\left(t_\text{flow,min} +
(t_\text{flow,design}-t_\text{flow,min})\cdot r^{1/n},\ t_\text{flow,min},\
t_\text{flow,design}\right), \\
r &= \max\!\left(0,\ \frac{t_\text{room}-t_\text{amb}}{t_\text{room}-t_\text{amb,design}}\right)
\end{align*}
mit den Vorgaben $t_\text{amb,design}=-12\gC$, $t_\text{flow,design}=110\gC$,
$t_\text{flow,min}=70\gC$, $t_\text{room}=20\gC$, $n=1$ (linear; $\approx$1{,}3
entspricht der Heizkörper-Charakteristik). Zwei \textbf{Presets} machen die
Temperaturabsenkungs-Geschichte mit einem Klick sichtbar:

\begin{itemize}[nosep]
  \item \ui{3G} --- \texttt{t\_flow\_design=110}, \texttt{t\_flow\_min=70}
    (ein klassisches 110/60-System);
  \item \ui{4G} --- \texttt{t\_flow\_design=70}, \texttt{t\_flow\_min=65}
    (ein Niedertemperatur-70/40-System).
\end{itemize}

Ein Netz ohne Heizkurven-Block läuft bei konstanter Vorlauftemperatur (die
Heizkurve ist dann \texttt{null} im Frame --- mehrere Referenznetze wie DESTEST
tun das). Dieselben Auslegungspunkte $t_\text{room}$/$t_\text{amb,design}$
verankern auch den Gradstunden-Faktor der Nachfrageskalierung
(Kapitel~\ref{ch:wetter}) --- die Physik bleibt konsistent, wenn der Betreiber
am Wetterregler zieht.

\begin{beispiel}[Wie sich die Verlustquote zwischen 3G und 4G ändert]
Auf dem Appendix-A-Netz bei $t_\text{amb}=0\gC$ macht ein Umschalten von 3G auf
4G beide gegenläufigen Effekte quantitativ:

{\small
\begin{center}
\begin{tabular}{@{}lrrrr@{}}
\toprule
Preset & $t_\text{flow}$ & Verluste & Massenstrom Zentrale & Pumpe $P_\text{el}$ \\
\midrule
3G (110/70) & 95{,}0\gC & 31{,}10~kW (16{,}3~\%) & 1{,}152~kg/s & 0{,}34~kW \\
4G (70/65)  & 68{,}1\gC & 21{,}35~kW (11{,}8~\%) & 2{,}199~kg/s & 0{,}64~kW \\
\bottomrule
\end{tabular}
\end{center}}

Die Verluste sinken um 31~\%, weil ein kühleres Vorlaufrohr weniger an das
Erdreich verliert. Aber: Für eine feste Nachfrage $Q=\dot m\,c_p\,\Delta T$
schrumpft bei kühlerem Vorlauf (und nicht mitsinkendem Rücklauf) das $\Delta T$,
also muss der Massenstrom steigen --- hier $\times 1{,}91$ --- und der
Pumpenstrom (mit $P_\text{hyd}=\dot V\cdot\Delta p$, quadratisch im Fluss)
nahezu verdoppelt sich. Das ist das \emph{Low-$\Delta T$-Syndrom}
(Niedertemperatur-Syndrom), direkt sichtbar in der Geschwindigkeits-Ebene der
Karte und in den Übersichts-Kennzahlen.
\end{beispiel}

\section{Die Schlechtpunktregelung}
\label{sec:schlechtpunkt}

Die Förderhöhe der Zentralen-Pumpe (\texttt{plift\_bar}) wird geregelt, um den
\emph{kritischen} Abnehmer versorgt zu halten. Der Regler macht \textbf{einen
begrenzten Proportionalschritt pro Tick, nach der Lösung} --- bewusst keine
innere Iterationsschleife:
\[
\texttt{plift} \mathrel{+}= \text{clamp}\big(K\cdot(\Delta p_\text{Soll}-\Delta p_\text{worst}),\ \pm\text{max\_step}\big)
\]
mit $K=0{,}5$, $\text{max\_step}=0{,}05$~bar/Tick, Sollwert-Band 0{,}3--2{,}0~bar
(Vorgabe 0{,}7~bar) und absoluter Förderhöhen-Grenze 0{,}1--10~bar. Die
aktualisierte Förderhöhe wirkt erst auf die Lösung des \emph{nächsten} Ticks ---
die Pumpe reagiert also \emph{über Minuten} wie eine echte drehzahlgeregelte
Pumpe, und die Engine bleibt bei einer Lösung pro Tick ($O(1)$, keine
Fixpunkt-Schleife). Der Modus \ui{fest} (\glqq ungeregelte Pumpe\grqq) hält die
Förderhöhe, wo der Nutzer sie gesetzt hat --- der Lehr-Kontrast, der die
Überdimensionierung einer Konstant-Förderhöhen-Pumpe zeigt. Dass der Regler
\emph{nur} die Gemessen-Ebene liest, ist der Kern der Blindflug-Semantik
(Abschnitt~\ref{sec:blindflug}).

\begin{beispiel}[Was passiert mit dem Schlechtpunkt-$\Delta p$, wenn morgens alle Stationen aufheizen?]
Fährt die morgendliche Nachfrage hoch, steigen die Massenströme und damit die
Druckverluste im Netz; der $\Delta p$ am ungünstigsten Abnehmer fällt. Der
Regler sieht das (über den bemessenen Schlechtpunkt) und hebt die Förderhöhe ---
aber nur um 0{,}05~bar je Tick. In der M4-Abnahme trieb er den beobachteten
Schlechtpunkt von 1{,}953~bar streng monoton auf den Sollwert 0{,}7~bar,
Setzung nach $\approx$26~Ticks bei 0{,}703~bar, dabei fiel der Pumpenstrom von
0{,}39 auf 0{,}14~kW. Mit der Vorlage \ui{Schlüsselstellen} lässt sich zusehen,
wie der Regler die Morgenrampe über Minuten nachfährt --- und wie das
\texttt{blind\_spot}-Flag flackert, wenn der wahre Schlechtpunkt wandert
(Abschnitt~\ref{sec:blindflug}). Ein überhöhter Sollwert kostet messbar
Pumpenstrom; der Modus \ui{fest} zeigt den Unterschied.
\end{beispiel}

% ===================================================================
\chapter{Anlagen}
\label{ch:anlagen}

Alle Anlagen lassen sich zur Laufzeit über das Element-Kontextmenü
(Rechtsklick $\rightarrow$ \glqq Hier platzieren\grqq) setzen und über den
angehefteten Abschnitt (Strg-Klick) konfigurieren. Sie wirken sofort im
laufenden Netz; ein einzelner nicht konvergierter Schritt heilt sich im nächsten
Takt selbst.

\section{Erzeuger}
\label{sec:erzeuger}

pandapipes weiß nichts von Kesseln oder Wärmepumpen --- der Slack ist nur eine
Druck-/Temperaturquelle. Die \emph{Erzeugerart} ist ein Dispatch-Modell
\emph{über} dem Slack, das aus der gelösten Einspeisung und den Live-Temperaturen
pro Tick die Strom-/Brennstoffseite berechnet:

{\small
\begin{center}
\begin{tabular}{@{}lll@{}}
\toprule
Art & Formel & Anmerkung \\
\midrule
Kessel (boiler) & $p_\text{fuel}=q/\eta$ & $\eta$ Vorgabe 0{,}95 \\
BHKW (chp)      & $p_\text{el}=\text{pq}\cdot q$ & wärmegeführt, pq Vorgabe 0{,}5 \\
Wärmepumpe (heat\_pump) & $\text{COP}=\eta_g\cdot T_\text{hot}/(T_\text{hot}-T_\text{cold})$ & $p_\text{el}=q/\text{COP}$ \\
\bottomrule
\end{tabular}
\end{center}}

$T_\text{hot}$ ist die \emph{live} Vorlauftemperatur der Zentrale, $T_\text{cold}$
die live Quelltemperatur ($t_\text{amb}$ für eine Luft-, $t_\text{ground}$ für
eine Erdreich-Wärmepumpe); der Gütegrad $\eta_g$ (Vorgabe 0{,}5) skaliert den
idealen Carnot-COP, der Temperaturhub ist bei 1~K abgefangen. Die Kennzahlen
werden jeden Tick neu gerechnet und erscheinen im \texttt{producers[]}-Block
neben \texttt{q\_kw}, \texttt{t\_flow\_c}, \texttt{plift\_bar} und
\texttt{pump\_el\_kw}.

\begin{beispiel}[Die Wärmepumpe belohnt niedrige Netztemperaturen]
$T_\text{hot}$ ist die live Vorlauftemperatur --- senkt die Heizkurve den
Vorlauf, steigt der COP und die berichtete elektrische Leistung sinkt. Das ist
die Temperaturabsenkungs-Geschichte, quantitativ: Das Absenken der
Netztemperaturen hebt den COP um grob 2--3~\% je Kelvin. Ein 4G-Netz mit einer
Erdreich-Wärmepumpe fährt deshalb sichtbar sparsamer als dasselbe Netz auf 3G ---
genau das vergleicht das mitgelieferte Szenario \glqq Demo-Dorf 4G Vergleich\grqq
(Kapitel~\ref{ch:szenarien}), das eine Erdreich-Wärmepumpe fährt.
\end{beispiel}

\section{Pufferspeicher}
\label{sec:speicher}

Ein Pufferspeicher (Vorgabe 100~kWh / 50~kW) überbrückt das Vor-/Rücklaufpaar
eines Knotens mit \emph{zwei} Zweigen, von denen pro Tick genau einer aktiv ist:

\begin{description}[nosep]
  \item[Laden] ein \texttt{heat\_consumer} mit \texttt{qext\_w} +
    \texttt{treturn\_k}~=~Speicher-Untertemperatur (Vorgabe 45\gC): zieht
    Ladeleistung aus dem Netz.
  \item[Entladen] eine \texttt{circ\_pump\_const\_mass\_flow} (druckfrei,
    \texttt{type="t"}) drückt einen festen Massenstrom bei der
    Speicher-Obertemperatur (Vorgabe 80\gC) in den Vorlauf. Die druckfreie
    Variante ist zwingend --- eine druckfixierende Pumpe würde gegen das
    Druckfeld des Slacks kämpfen und alle Retry-Stufen zum Divergieren bringen.
  \item[Bereitschaft] der Ladezweig fällt auf das Nullfluss-Bodenpaar
    (0{,}02~kg/s + 100~W), die Entladepumpe geht außer Betrieb --- die
    Stummelrohre erreichen nie Nullfluss.
\end{description}

Der Ladestand (SoC) wird aus den \emph{gelösten} Ergebnissen integriert (Laden
über die realisierte \texttt{qext\_w}, Entladen über $\dot m\,c_p\,\Delta T$),
mit Einweg-Wirkungsgrad $\eta=0{,}95$ (Umlauf $\eta^2\approx0{,}90$); die
Bereitschaft ist ein stehender Verlust, der SoC bewegt sich dann nicht.

\begin{beispiel}[Warum ein voller Speicher von selbst aufhört zu laden]
\texttt{desired\_state} wendet Leistungs- \emph{und} Kapazitätsgrenzen an,
bevor der aktive Zweig feststeht: Eine Ladeanforderung wird durch den
verbleibenden Raum gedeckelt, eine Entladeanforderung durch die verfügbare
Ladung; unter 1~W idlet der Speicher von selbst. Ein voller Speicher stellt das
Laden also selbsttätig ein, ein fast leerer klingt beim Entladen geometrisch ab
und geht dann in Bereitschaft. Zum Kappen der morgendlichen Wärmespitze lädt man
den Speicher nachts (niedrige Last, günstige Erzeugung) und entlädt ihn zur
Morgenrampe --- der SoC-Balken im angehefteten Abschnitt zeigt das live.
\end{beispiel}

\section{Einspeiser und Bypass}

\paragraph{Einspeiser} (Solarthermie/Abwärme, Vorgabe 20~kW) sind
Wärmetauscher: eine feste Einspeisung als negatives \texttt{qext\_w}. Setzt man
einen 20-kW-Einspeiser, fährt die Zentrale entsprechend zurück (verifiziert:
187{,}2~$\rightarrow$~173{,}8~kW, Bilanz 0{,}07~\%). Die Plattform bleibt ein
Zweileiternetz mit \emph{einem} Druckhalter --- weitere Erzeuger speisen als
Wärmetauscher oder Massenstrompumpen ein, nie als zweiter Druck-Slack (der
zweite Slack quittiert mit HTTP~409).

\paragraph{Bypass} (Netzschluss) ist das kanonische Massenstrom-Paar: ein
winziger fester Massenstrom (0{,}02~kg/s) plus eine kleine Bereitschaftswärme
(100~W), gesetzt, um ein Endstrang durchströmt und warm zu halten.

\begin{beispiel}[Der Bypass hält den Endstrang warm]
Ein Trassenknoten mit Grad~1 (ein Endstrang ohne Abnehmer) trägt ohne Anschluss
Nullfluss --- und ein Nullfluss-Ast ist in der Wärme-Matrix von pandapipes
singulär. Der Bypass löst das: sein fester Massenstrom trägt keine
Singularität. Seine Bereitschaftswärme parkt im Warmwasser-Slot, damit der
Wetter-Override sie nie umskaliert --- die Aufgabe eines Bypasses ist
Temperaturhaltung, nicht wetterfolgende Nachfrage. Das reale Ortsnetz
Schutterwald setzt genau diesen kanonischen Bypass an seine abnehmerlosen
Zweige.
\end{beispiel}

\section{Die Nullfluss-Schwelle}
\label{sec:nullfluss}

Ein \texttt{heat\_consumer} mit null \texttt{qext\_w} trägt null Massenstrom,
und die Nahe-Null-Last ist genau das Regime, in dem die
\texttt{qext\_w}+\texttt{treturn\_k}-Regelung divergiert (es gibt keinen Fluss,
dessen Rücklauftemperatur man regeln könnte). Die Schwelle
\api{RTHEATFLOW\_MIN\_QEXT\_W} (Vorgabe 500~W) hebt die Gesamtnachfrage jedes
\emph{temperaturgeregelten} Abnehmers auf diesen Boden. Entscheidend: Die
Schwelle ist \emph{zeilenweise und maskengerecht} --- Massenstrom-Paar-Zeilen
(Bypässe, festfließende Abnehmer) sind ausgenommen, ein 100-W-Bypass bleibt
100~W. Der Vorgabewert 500~W ist ein bewusst großzügiger Löser-Schutz für den
interaktiven Betrieb; die Validierungsläufe senken ihn unter das reale Minimum
des jeweiligen Netzes (DESTEST 62~W).

% ===================================================================
\chapter{Der Wetter-Regler}
\label{ch:wetter}

Eine Wetter-Zeitreihe treibt alles: Die Außentemperatur speist die
Raumheizlast und die Heizkurve, die Erdreich-Temperatur speist \texttt{text\_k}
jedes Rohres. Der Abschnitt \ui{Wetter} in der Seitenleiste zieht die
Außentemperatur live ($-30$ bis $+45\gC$, per \api{PUT /weather/override},
150-ms-entprellt); Loslassen kehrt zum Wetterprofil zurück.

\section{Die Gradstunden-Skalierung}

Der Mechanismus ist der Gradstunden-Faktor
$f(T)=\max(0,T_\text{room}-T)/(T_\text{room}-T_\text{design})$, angewandt
\emph{nur} auf die Raumheizung --- das Warmwasser (DHW) bleibt \textbf{unberührt},
weil die Warmwassernachfrage dem Wetter nicht folgt. Zwei Regime, abgesichert
durch $\varepsilon=0{,}05$:

\begin{itemize}[nosep]
  \item \textbf{Normal:} Das Profil hat nennenswerte Raumheizung
    ($f_\text{Profil}\geq\varepsilon$) $\rightarrow$ skaliere es mit dem
    Verhältnis $f(T_\text{override})/f(T_\text{Profil})$ (die Tagesform bleibt
    erhalten).
  \item \textbf{Sommer:} Die Raumheizung des Profils ist $\approx$0
    ($f_\text{Profil}<\varepsilon$) $\rightarrow$ es gibt kein Verhältnis zum
    Skalieren, also baue die Nachfrage aus der Auslegungslast neu auf,
    $q_\text{design}\cdot f(T_\text{override})$.
\end{itemize}

Der Override setzt nur die effektive Außentemperatur; die Erdreich-Temperatur
wird nie überschrieben (eine langsame, saisonale Größe). Weil dieselben
Auslegungspunkte auch die Heizkurve verankern, reagieren Vorlauftemperatur
\emph{und} Nachfrage konsistent auf den Regler.

\begin{beispiel}[Einen Sommertag in den Winter ziehen]
Startet man an einem Sommertag (Raumheizung nahe null) und zieht den
Wetterregler auf $-10\gC$, greift das Sommer-Regime: Es gibt keine Tagesform
zum Skalieren, also rekonstruiert die Plattform die Heizlast aus der
Auslegungslast mal Gradstundenfaktor --- man bekommt eine physikalisch
plausible Winterlast aus einem Nahe-Null-Profil. So lassen sich Kälteeinbruch
und Sommerbetrieb im Zeitraffer durchspielen, ohne das Netz zu wechseln.
\end{beispiel}

\begin{beispiel}[Warum die \emph{relative} Verlustquote im Sommer hochspringt]
Die Wärmeverluste eines Netzes sind fast nachfrageunabhängig (die stehenden
Verluste des Rohrs bleiben, ob viel oder wenig Wärme fließt). Zieht man den
Wetterregler in den Sommer, bricht die Nachfrage ein --- die \emph{absoluten}
Verluste sinken kaum, die \emph{relative} Verlustquote (Verlust/Einspeisung)
springt sichtbar hoch. Am realen Schutterwald reicht dieselbe Physik von
$\approx$7{,}5~\% im tiefen Winter bis $\approx$45~\% im Sommer. Das ist derselbe
Effekt wie die Abhängigkeit von der Wärmedichte --- nur in der Zeit statt im
Raum.
\end{beispiel}

% ===================================================================
\chapter{Referenznetze und validierte Physik}
\label{ch:referenznetze}

Neben den Lehrnetzen enthält der Katalog vier Netze aus realen bzw.\
publizierten Quellen. Jedes Netzverzeichnis unter \api{data/networks/<id>/}
trägt eine \texttt{DATASET.md} mit Herkunft, Lizenz,
Konvertierungsentscheidungen und dem vollständigen Validierungsprotokoll ---
und alle hier genannten Zahlen sind \emph{live nachgerechnete} Werte, gegen
\emph{externe} Wahrheit geprüft, mit vorab (nicht nachträglich) begründeten
Toleranzbändern.

\section{Der Katalog}

{\small
\begin{center}
\begin{tabular}{@{}llrrl@{}}
\toprule
id & Charakter & Knoten & Trasse km & Rolle \\
\midrule
\texttt{demo\_dorf}   & rural     & 15  & 0{,}953 & Karten-Demo (Standardnetz) \\
\texttt{appendix\_a}  & abstract  & 4   & 1{,}0   & bekannte Antwort-Referenz \\
\texttt{schutterwald} & rural     & 208 & 2{,}626 & reales Ortsnetz (Vorzeigenetz) \\
\texttt{destest\_16}  & benchmark & 25  & 0{,}39  & \textbf{validiert} (Tabelle 6) \\
\texttt{destest\_8}   & benchmark & 13  & 0{,}216 & Katalogbeispiel \\
\texttt{destest\_32}  & benchmark & 49  & 0{,}792 & Katalogbeispiel \\
\texttt{verbier}      & meshed    & 676 & 12{,}05 & vermascht, Messdaten \\
\bottomrule
\end{tabular}
\end{center}}

\paragraph{Die Lehrnetze.} \texttt{demo\_dorf} ist ein handgefertigtes Straßendorf
auf realer Geographie (nördlich von Karlsruhe): eine \glqq Heizzentrale\grqq,
11~Abnehmer (6$\times$EFH\_ALT, 3$\times$EFH\_SAN, 2$\times$MFH), 182~kW
Auslegung, 3G-Heizkurve --- es existiert, damit die Karte etwas Reales zeichnen
kann (lineare Wärmedichte bewusst unter der Faustregel, siehe
Abschnitt~\ref{sec:netzstudio}). \texttt{appendix\_a} ist die abstrakte
Vier-Knoten-Referenzschleife, deren Abnehmer A/B/C \emph{alle drei}
Partnertypen ausüben (Rücklauf/$\Delta T$/Massenstrom); ihr Slack trägt
\emph{keine} Heizkurve, damit die Antwort deterministisch ist --- die Testsuite
pinnt an ihr die bekannten Zahlen (A~$\dot m$~0{,}674357, C~$\dot m$~0{,}25,
Verluste 27{,}258~kW, Einspeisung 187{,}182~kW, Bilanz $-0{,}04~\%$).

\section{DESTEST --- der validierte Cross-Tool-Benchmark}

Die IBPSA-\glqq DESTEST\grqq-Übungen sind ein Cross-Tool-Benchmark: mehrere
unabhängige Werkzeuge modellieren dasselbe kleine Fernwärmenetz und
veröffentlichen ihre Ergebnisse, sodass die Streuung zwischen den Werkzeugen die
Toleranz ist. \texttt{destest\_16} (16~Gebäude, Netz CE\_1, aus der
Fallbeschreibung Tabelle~6 gebaut) wird gegen zwei Übungen validiert. Der
Solver läuft mit \texttt{min\_qext\_w~=~62}~W --- bei $\Delta T$~30~K ist das die
1{,}77-kg/h-Mindestumgehung von DESTEST, der Nullfluss-Boden \emph{ist} die
Bypass-Regel des Falls.

\paragraph{CE 0 --- eingeschwungener Zustand} (alle 16~Stationen exakt auf
Spitzenlast, gegen 6~Werkzeuge):

{\small
\begin{center}
\begin{tabular}{@{}llll@{}}
\toprule
Kennzahl & rtheatflow & publiziert (6 Tools) & Urteil \\
\midrule
Massenstrom Zentrale & 8876{,}5~kg/h & 8847{,}9--8870{,}4 & $+0{,}07~\%$ über Max \\
Rücklauftemperatur   & 39{,}48\gC    & 39{,}46--39{,}85   & \textbf{im Band} \\
Wärme gesamt         & 314{,}8~kW    & 308{,}2--314{,}3   & $+0{,}16~\%$ über Max \\
Vorlauf an SD1       & 69{,}454\gC   & 69{,}43--69{,}48   & \textbf{im Band} \\
\bottomrule
\end{tabular}
\end{center}}

Die beiden Kennzahlen knapp über dem publizierten Maximum (Massenstrom und
Wärme) sind genau die, die der Fluid-Eigenschaftsunterschied bewegen soll:
pandapipes' temperaturabhängiges Wasser läuft mit $c_p\approx4179$--$4187$~J/(kg$\cdot$K)
gegen DESTESTs feste 4180, und ein $\Delta T$-geregelter Massenstrom skaliert
invers mit $c_p$. Die Vorlauftemperatur an SD1 --- reine Verlustphysik, vom
Fluid-Konstanten unabhängig --- landet \emph{im} Band: der eigentliche
Prüfstein, dass das mehrschichtige Wandverlustmodell die publizierten
Rohrverluste reproduziert.

\paragraph{CE 1 --- die Sieben-Tage-Woche} (1008~Lösungen auf dem 10-min-Raster,
gegen 3~dynamische Werkzeuge):

{\small
\begin{center}
\begin{tabular}{@{}lrrrr@{}}
\toprule
Kennzahl & rtheatflow & AixLib & Buildings & IBPSA \\
\midrule
Wärme-Einspeisung & 14{,}506~MWh & 14{,}452 & 14{,}324 & 14{,}357 \\
Netzverluste      & 606{,}6~kWh  & 535{,}2  & 543{,}8  & 535{,}7 \\
\bottomrule
\end{tabular}
\end{center}}

\begin{beispiel}[Die ehrliche Nacht-Zuordnung der Verluste]
Der Verlust-Gesamtwert liegt $\approx$12~\% über dem publizierten Maximum ---
und das ist die lehrreichste ehrliche Zuordnung der Plattform. Die Woche
zerfällt sauber: Die \emph{belasteten} Ticks (Nachfrage $>$5~kW) tragen
\textbf{534~kWh} Verlust --- \emph{innerhalb} des publizierten Bandes
535--544~kWh. Unter Last stimmt die Verlustphysik mit den dynamischen
Rohrwerkzeugen überein. Der \textbf{Überschuss von $+72{,}5$~kWh} sitzt komplett
in den \textbf{401 Nahe-Null-Nachtticks}: Dort kühlen die publizierten
dynamischen Rohre \emph{ab}, während ein quasistatischer eingeschwungener
Zustand die Rohreintritts-Region auf Vorlauftemperatur hält und deshalb den
stehenden Verlust überzählt. Der Überschuss ist nicht frei --- er ist \emph{vor}
der Messung durch die Enthalpie der DESTEST-Mindestumgehung nach oben begrenzt
($16\times1{,}77$~kg/h$\times c_p\times60$~K $\approx$~2~kW über 401~Ticks
$\Rightarrow\leq$~134~kWh). Das Band ist also asymmetrisch und \emph{hergeleitet},
nicht getunt.
\end{beispiel}

\section{Schutterwald --- Plausibilität am realen Ort}

Schutterwald ist ein realer Ort (Baden) ohne publizierte Cross-Tool-Antwort,
deshalb wird es auf \emph{Plausibilität} geprüft: geschlossene Energiebilanz,
Geschwindigkeiten im Band, Jahres-Verlustquote im Literaturbereich. 208~Trassenknoten,
207~Trassen, 44~Übergabestationen aus dem realen Gasnetz abgeleitet
(50-m-Anschlussradius), 2{,}626~km reale WGS84-Trasse, 1904~kW Auslegung, LHD
1{,}34~MWh/(m$\cdot$a), 3G-Kurve. Das Karten-Vorzeigenetz.

{\small
\begin{center}
\begin{tabular}{@{}lll@{}}
\toprule
Metrik & Wert & Prüfung \\
\midrule
Konvergenz            & alle 96 Schritte Tier 1 & \texttt{solver\_status="ok"} \\
worst-case Bilanzfehler & 0{,}075~\% der Einspeisung & $\leq$1~\% \\
Spitzengeschwindigkeit & 0{,}98~m/s & $<$1{,}3~m/s \\
Wintertag-Verlustquote & \textbf{7{,}5~\%} der Einspeisung & $0{,}04<r<0{,}12$ \\
\bottomrule
\end{tabular}
\end{center}}

Die Verluste sind nahezu nachfrageunabhängig ($\approx$61~kW etwa konstant), also
annualisiert die instantane Winterquote von 7{,}5~\% höher: 61~kW$\times$8760~h
$\approx$537~MWh Verlust gegen $\approx$3{,}52~GWh Nachfrage $\Rightarrow$
$\approx$\textbf{13~\%} Jahresverlustquote --- im 8--20-\%-Plausibilitätsband für
ein 3G-Netz dieser Dichte. Die Unterscheidung instantan/annualisiert ist der
Punkt: Ein einzelner kalter Tag unterschätzt die Jahresverluste.

\section{Verbier --- gegen reale Messdaten}

Verbier (OpenDHN, Zenodo DOI 10.5281/zenodo.10793816, CC~BY~4.0) ist das einzige
Referenznetz, das gegen \emph{reale eingeschwungene Messdaten} validiert wird ---
ein reales vermaschtes Alpennetz: 676~Trassenknoten, 681~Trassen je Seite
(12{,}05~km, \textbf{6~unabhängige Schleifen}), zwei Heizzentralen, 150~Übergabestationen
mit gemessener Leistung und gemessenem Rücklauf. Die gemessenen
\emph{Vorlauf}temperaturen sind \emph{nie} Eingabe --- sie sind das
Validierungsziel. Der vermaschte 1352-Knoten-Netz mit druckfreier zweiter Pumpe
scheitert an Tier~1 und konvergiert erst auf \textbf{Tier~2} ($\alpha=0{,}5$);
kalt $\approx$28~s, warm $\approx$5~s je Lösung --- deshalb ist Verbier als
Batch-/Analysenetz dokumentiert, nicht als Echtzeitnetz.

{\small
\begin{center}
\begin{tabular}{@{}llll@{}}
\toprule
Kennzahl & rtheatflow & gemessen & Delta \\
\midrule
Einspeisung HS0+HS1 & 3567~kW    & 3590~kW    & $-0{,}6~\%$ \\
HS0 Massenstrom     & 17{,}83~kg/s & 17{,}76~kg/s & $+0{,}4~\%$ \\
HS0 Rücklauf        & 44{,}10\gC  & 43{,}52\gC  & $+0{,}58$~K \\
Netzverluste        & 297~kW     & $\approx$317~kW & $-6~\%$ \\
\bottomrule
\end{tabular}
\end{center}}

\begin{beispiel}[Median 0{,}83~K über 150 Stationen]
Der Kern der Validierung sind die 150~Übergabestations-Vorlauftemperaturen ---
aus den Eingaben \emph{herausgehalten}: Median $|\Delta T|$~\textbf{0{,}83~K}
(Toleranz $\leq$2~K), Mittel 1{,}29~K, p90 2{,}68~K (Toleranz $\leq$5~K), Max
12{,}6~K (Toleranz $\leq$20~K). Ein medianer Absolutfehler von 0{,}83~K über
150~unabhängig gemessene Stationen, mit den Vorlauftemperaturen aus den
Eingaben gehalten, ist das stärkste Validierungsergebnis der Plattform: Die
gekoppelte hydraulisch-thermische Lösung rekonstruiert das reale Temperaturfeld
im Median bis auf Messrauschen. Fünf Stationen dominieren die Ausreißerspitze
(z.\,B.\ S96 mit $+12{,}6$~K) --- in der \texttt{DATASET.md} als
Sensor-Mittelung, uniforme 5\gC-Randtemperatur und lokale Fluss-Aufteilung im
Netz benannt und analysiert. Der \emph{Median}, nicht das Maximum, ist die
Aussage, weil dies reale Messdaten sind.
\end{beispiel}

\section{Akzeptierte Modellunterschiede}

Jedes Band oben ist durch dieselbe kurze Liste bewusster, dokumentierter
Entscheidungen geweitet --- keine ist ein Fehler, jede ist in der jeweiligen
\texttt{DATASET.md} begründet: temperaturabhängiges Wasser vs.\ feste
Eigenschaften ($\pm0{,}3$--$0{,}5~\%$ auf $\Delta T$-geregelte Größen);
quasistatisches Wärmemodell vs.\ dynamische Rohre (der CE-1-Nachtüberschuss,
oben hergeleitet begrenzt); angenommene Randtemperatur, wo die Quelle schweigt
(Verbier 5\gC, $\pm$5~K~$\Rightarrow$~$\mp$7~\% Verluste --- größer als das
$-6~\%$-Delta, das damit \emph{kein} Modellfehler ist); reale Messdaten
$\Rightarrow$ bewusst großzügige Toleranzen (der Median zählt, nicht das Maximum).

% ===================================================================
\chapter{Szenarien}
\label{ch:szenarien}

Das Menü \ui{Datei $\rightarrow$ Szenario speichern…} legt den kompletten
Aufbau als \textbf{Rezept} ab --- nicht als Schnappschuss: die Netz-id, die
deterministische Lastpolitik, die zur Laufzeit platzierten Anlagen (als
Konfiguration, nicht als Zustand), das Op-Protokoll der Abnehmer/Bypässe, die
Heizkurven-/$\Delta p$-/Erzeuger-Konfiguration, der Wetter-Override und die
Engine-Uhr. Weil jede Zutat deterministisch ist (geseedete Profile,
protokollierte Ops), reproduziert das Nachspielen den Aufbau exakt, und die
Datei bleibt lesbar --- bewusst handeditierbares JSON unter
\api{data/scenarios/}. Das Laden ist \emph{tolerant je Eintrag}: Passt eine Op
nach einem Netzwechsel nicht mehr, wird sie mit einer Warnung übersprungen statt
mit einem halben 500.

Mitgeliefert sind zwei Szenarien mit \textbf{gleichem Seed 42}:
\texttt{demo-dorf-3g-winter} und \texttt{demo-dorf-4g-vergleich}.

\begin{beispiel}[Direkter Vergleich der Temperaturniveaus]
Weil beide Szenarien denselben Seed tragen, sind ihre Lastprofile identisch ---
der einzige Unterschied ist das Temperaturniveau (und die 4G-Variante fährt eine
Erdreich-Wärmepumpe). Lädt man abwechselnd beide, sieht man am selben Netz und
denselben Lasten direkt, wie die Verlustquote fällt, der Massenstrom (und
Pumpenstrom) steigt und der Wärmepumpen-COP im 4G-Fall besser wird ---
Kapitel~\ref{sec:heizkurve} und~\ref{sec:erzeuger}, aber als greifbarer
Seite-an-Seite-Vergleich statt als Tabelle.
\end{beispiel}

% ===================================================================
\chapter{Daten exportieren}
\label{ch:export}

\section{Aufzeichnung}

\ui{Datei $\rightarrow$ Aufzeichnung starten}: Jeder veröffentlichte Schritt
wandert über eine nicht-blockierende Senke (eigener Schreib-Thread + Queue, die
Engine blockiert nie) in ein CSV-Paket unter \api{data/recordings/<id>/}. Das
Paket enthält je eine CSV für \texttt{summary}, \texttt{junctions},
\texttt{pipes}, \texttt{consumers}, \texttt{producers}, \texttt{storages},
\texttt{controls} und die Messwerte, dazu eine \texttt{metadata.json} als
Reproduktions-Rezept. Weil die Senke den \emph{projizierten} Frame konsumiert,
gilt der strikte Modus auch für die CSVs: Was nie auf den Draht kommt, kommt nie
in die Dateien. Beenden, dann als ZIP herunterladen (\ui{Aufzeichnungen}-Flyout).

\section{Tage-Export (Bulk)}

\ui{Datei $\rightarrow$ Tage exportieren…} simuliert ganze Tage des aktuellen
Aufbaus offline so schnell wie möglich: Die Engine wird kurz geparkt, der
Simulator tief kopiert, auf einen deterministischen Mitternacht-Start
normalisiert (\texttt{prepare\_replay}: Kaltstart, SoC~0, frische
Messfenster, geregelte Pumpe auf die Datei-Förderhöhe gelöst) und dann die Tage
zurück in ein privates Paket abgespielt. Das Paket ist \textbf{byte-kompatibel}
zu einer Live-Aufzeichnung.

\begin{beispiel}[Was \glqq byte-kompatibel\grqq{} bedeutet]
Ein verkürzter Tag wird einmal live über den Veröffentlichungspfad
aufgezeichnet, dann aus einer tiefen Kopie des \emph{driftenden} Simulators
(Warmstart fortgeschritten, $\Delta p$-Pumpe bewegt, Fenster gefüllt) exportiert
--- und jede CSV ist \emph{byte-identisch}, nachdem genau zwei dokumentierte
volatile Spalten maskiert sind: \texttt{timestamp} (Wanduhr) und
\texttt{solve\_ms} (Maschinen-Timing). Das ist der Beweis, dass der Offline-Export
dieselbe Physik rechnet wie der Live-Betrieb --- nur schneller. Die
Schätz-Ebene wird bewusst \emph{nicht} aufgezeichnet (ihre Kadenz ist
Maschinen-Timing, kein Physik-Takt; sie ist aus den aufgezeichneten Messwerten
neu berechenbar).
\end{beispiel}

\section{Der experimentelle Transient-Modus}
\label{sec:transient}

\api{RTHEATFLOW\_TRANSIENT=true} rechnet \emph{nur den Export} mit thermischer
Trägheit der Wassersäule (nicht Rohrwand oder Erdreich): Der Solver stellt zwei
transiente Stufen (mit explizitem \texttt{dt} und monotonem
\texttt{simulation\_time\_step}) vor die quasistatische Ladder. Das Verketten
pro Schritt ist genau das, was \texttt{run\_timeseries} intern tut, und wurde
\emph{bit-identisch} nachgewiesen, bevor darauf gebaut wird. Scheitert eine
transiente Stufe, fällt der Schritt automatisch auf quasistatisch zurück und das
Paket wird markiert (\texttt{transient\_fallback}). Der \emph{Live}-Betrieb
bleibt dagegen \emph{immer} quasistatisch --- der Transient-Modus ist
ausschließlich der Offline-Export, nie die Echtzeitschleife.

\begin{beispiel}[Trägheit sichtbar machen]
Hebt man im Transient-Export die Vorlauftemperatur der Zentrale an, sieht das
ferne Netzende die neue Temperatur \emph{nicht} sofort: Sie zerfällt über
Schritte ($358{,}05\rightarrow357{,}57$~K) statt des quasistatischen sofortigen
Sprungs ($351{,}8$~K) --- die Transportverzögerung, über 1~K, per Test gepinnt.
Für sinnvolle Ergebnisse wählt man die Rohr-\texttt{sections} so, dass ein
Abschnitt etwa $v\cdot dt$ lang ist.
\end{beispiel}

% ===================================================================
\chapter{Langzeit-Visualisierung mit Grafana}
\label{ch:grafana}

\texttt{docker compose up} startet den vollen Visualisierungsstack:

\begin{enumerate}[nosep]
  \item \textbf{Collector} --- pollt \api{GET /state} alle 0{,}5~s, entdoppelt
    auf \texttt{(day, step)} und schreibt wanduhr-gestempelte Punkte nach
    InfluxDB: \texttt{summary} (alle Kennzahlen + \texttt{converged} +
    Solverstatus + $\Delta p$-Soll/Ist + Wetter), \texttt{junction},
    \texttt{pipe}, \texttt{consumer}, \texttt{producer}, \texttt{storage}.
  \item \textbf{InfluxDB 2.7} --- Bucket \texttt{heatflow}, Organisation
    \texttt{rtheatflow}, 7~Tage Aufbewahrung (Port~8087).
  \item \textbf{Grafana 11} --- vollständig aus dem Repository provisioniert
    (Datenquelle + Dashboard-JSON), Port~3001, Zugang admin/admin.
\end{enumerate}

Die mitgelieferten Panels des Fernwärme-Dashboards: Vorlauf/Rücklauf der
Zentrale gegen den Heizkurven-Sollwert; Schlechtpunkt-$\Delta p$ gegen Sollwert
und Förderhöhe (mit 0{,}5-bar-Schwellenlinie); Erzeuger-Dispatch als Stapel;
Speicher-SoC; Verlustquote als Anzeige; Rechenzeit als Statwert; und der
Solverstatus als OK/DEGRADED/FAILED-Abbildung. Sieht Grafana nichts, wartet der
Collector, bis rtheatflow und InfluxDB gesund sind und der erste Schritt gelöst
ist --- Zeitbereich auf \glqq letzte 5~Minuten\grqq{} stellen.

% ===================================================================
\chapter{Konfiguration}
\label{ch:konfiguration}

Alle Einstellungen erfolgen über Umgebungsvariablen mit dem Präfix
\texttt{RTHEATFLOW\_} (eine \texttt{.env}-Datei wird gelesen; Vorlage:
\texttt{.env.example}). In der Tabelle ist das Präfix weggelassen.

{\small
\begin{longtable}{@{}p{0.33\textwidth}p{0.16\textwidth}p{0.42\textwidth}@{}}
\toprule
Variable & Vorgabe & Bedeutung \\
\midrule
\endhead
\texttt{DATA\_DIR}              & \texttt{./data}     & Datenverzeichnis \\
\texttt{DEFAULT\_NETWORK}       & \texttt{demo\_dorf} & Netz beim Start \\
\texttt{STEP\_INTERVAL\_SECONDS}& \texttt{1.0}        & echte Sekunden je Simulationsminute (0{,}01\,…) \\
\texttt{STEPS\_PER\_DAY}        & \texttt{1440}       & Schritte pro Tag \\
\texttt{AUTOSTART}              & \texttt{true}       & Simulation beim Start anlaufen lassen \\
\texttt{HISTORY\_SIZE}          & \texttt{1440}       & Größe des Verlaufspuffers \\
\texttt{SOLVER\_ITER}           & \texttt{100}        & Basis-Iterationen der Retry-Ladder \\
\texttt{PUMP\_ETA}              & \texttt{0.7}        & elektr.\ Pumpenwirkungsgrad ($P_\text{el}=P_\text{hyd}/\eta$) \\
\texttt{DP\_MIN\_BAR}           & \texttt{0.5}        & Mindest-$\Delta p$ am kritischen Abnehmer [bar] \\
\texttt{RETURN\_TEMP\_MARGIN\_K}& \texttt{5}          & Bernstein-Schwelle Rücklaufverletzung [K] \\
\texttt{MIN\_QEXT\_W}           & \texttt{500}        & Nullfluss-Schutz: Abnehmer-Lastuntergrenze [W] \\
\texttt{EXPOSE\_GROUND\_TRUTH}  & \texttt{true}       & \texttt{false} = strikter Modus (Realität verlässt den Server nicht) \\
\texttt{TRANSIENT}              & \texttt{false}      & \texttt{true} = Offline-Export mit thermischer Trägheit (experimentell) \\
\texttt{RECORD}                 & \texttt{false}      & \texttt{true} = Daueraufzeichnung, ein Paket je Konfiguration \\
\texttt{HOST} / \texttt{PORT}   & \texttt{127.0.0.1} / \texttt{8001} & API-Bindung (Docker/LAN: \texttt{0.0.0.0}); Geschwisterschema \\
\texttt{CORS\_ORIGINS}          & \texttt{*}          & erlaubte Browser-Origins \\
\texttt{LOG\_LEVEL}             & \texttt{info}       & Log-Ausführlichkeit \\
\bottomrule
\end{longtable}}

\begin{beispiel}[.env für eine Prüfungsumgebung]
Doppeltes Tempo, strikte Beobachtbarkeit (Studierende können die Wahrheit auch
per API nicht abfragen), Zugriff aus dem Kursnetz:
\begin{lstlisting}
RTHEATFLOW_STEP_INTERVAL_SECONDS=0.5
RTHEATFLOW_EXPOSE_GROUND_TRUTH=false
RTHEATFLOW_HOST=0.0.0.0
\end{lstlisting}
\end{beispiel}

% ===================================================================
\appendix

\chapter{REST-/WebSocket-Schnittstelle}
\label{app:api}

Für Skripte, eigene Frontends oder Übungsaufgaben ist die komplette
Funktionalität auch per HTTP erreichbar (Standard:
\url{http://127.0.0.1:8001}, ohne Authentifizierung --- Lehrwerkzeug). Die
Plattform hat \textbf{61 Routen}; die \emph{vollständige} Referenz steht in
\api{docs/API.md} (aus dem Code generiert) und interaktiv in Swagger unter
\api{/docs}. Auswahl der wichtigsten Kategorien:

{\small
\begin{longtable}{@{}p{0.17\textwidth}p{0.34\textwidth}p{0.40\textwidth}@{}}
\toprule
Methode & Pfad & Beschreibung \\
\midrule
\endhead
GET & \api{/health}, \api{/status} & Diagnose, Engine-Zustand \\
GET & \api{/network}               & statische Topologie (inkl.\ Geokoordinaten) \\
GET & \api{/state}                 & letzter gelöster Schritt (404 vor der ersten Lösung) \\
GET & \api{/history?limit=N}       & Verlaufspuffer \\
WS  & \api{/ws}                     & ein projizierter \texttt{StepResult} je Schritt \\
GET & \api{/manual}                & dieses Handbuch (HTML; \texttt{?format=md} roh) \\
POST& \api{/control/start|pause|resume} & Simulation steuern \\
POST& \api{/control/seek|seekday|interval} & Zeit/Tag/Schrittdauer setzen \\
GET/PUT/DEL & \api{/weather\ldots}  & Wetter + Außentemperatur-Override \\
GET/POST & \api{/heatingcurve}, \api{/dpcontrol} & Heizkurve, Schlechtpunktregelung \\
GET/POST/DEL & \api{/producer\ldots}, \api{/storage\ldots} & Erzeuger, Speicher \\
POST/DEL & \api{/consumer\ldots}, \api{/bypass} & Abnehmer, Bypass \\
GET/POST/DEL & \api{/measurements\ldots} & Messstellen, Presets, Modus \\
GET/POST & \api{/estimation/config} & Beobachter-Richtlinie \\
GET/POST & \api{/networks\ldots}, \api{/config/apply} & Katalog, Import, Netz laden \\
POST & \api{/loadgen/assign}        & Lastpolitik-Vorschau \\
GET/POST/DEL & \api{/scenarios\ldots} & Szenarien speichern/laden/löschen \\
GET/POST & \api{/recording\ldots}, \api{/export\ldots} & Aufzeichnung, Bulk-Export \\
\bottomrule
\end{longtable}}

Fehlercode-Konventionen: \texttt{400} Validierung/Import, \texttt{404} fehlende
Ressource (oder \api{/state} vor der ersten Lösung), \texttt{409} Konflikte
(zweiter Druck-Slack, laufender Export), \texttt{422} semantische Grenzen
(Wetter-Override außerhalb $-30\dots45\gC$), \texttt{500} nur interne Fehler ---
\textbf{Solver-Nichtkonvergenz ist Daten (\texttt{converged=false}-Frames), nie
ein 500.}

\section{Das StepResult-Wireformat}

Das einzige Wireformat ist der projizierte \texttt{StepResult}: \api{/state},
die \api{/history}-Einträge und jede \api{/ws}-Nachricht teilen einen
\texttt{asdict()}+Projektions-Pfad. Im strikten Modus fehlen die vier
Wahrheits-Schlüssel (\texttt{junctions}, \texttt{pipes}, \texttt{consumers},
\texttt{summary}), \texttt{measurements}/\texttt{observed\_summary} und
\texttt{estimated} bleiben.

\begin{lstlisting}
$ curl http://127.0.0.1:8001/state
{
  "step": 440, "day": 0, "time_of_day": "07:20",
  "converged": true, "solver_status": "ok", "solve_ms": 24.1,
  "junctions": [...], "pipes": [...], "consumers": [...],
  "summary": { "q_feed_kw": ..., "q_loss_kw": ..., "loss_pct": ...,
               "dp_worst_bar": ..., "worst_consumer": "...",
               "t_flow_plant_c": ..., "balance_err_kw": ... },
  "producers": [...], "storages": [...],
  "weather": {...}, "controls": { "dp_control": { "blind_spot": false, ... } },
  "measurements": { "consumers": [...], "nodes": [...], "plant": {...} },
  "observed_summary": { "dp_worst_bar": ..., "n_metered": ... },
  "estimated": { "summary": {...}, "error": { "max_dt_return_k": ... } }
}
\end{lstlisting}

\begin{beispiel}[Drei typische Aufrufe]
Ein Netz mit Lastpolitik laden (entspricht \ui{Anwenden}):
\begin{lstlisting}
curl -X POST http://127.0.0.1:8001/config/apply \
  -H "Content-Type: application/json" \
  -d '{"network_id": "schutterwald",
       "loadgen": {"temperature_preset": "3G", "seed": 42}}'
\end{lstlisting}
Einen Wärmemengenzähler an Abnehmer 5 setzen und den Rücklauf abfragen:
\begin{lstlisting}
curl -X POST http://127.0.0.1:8001/measurements/consumer/5
\end{lstlisting}
Den Live-Stream anzapfen --- jede Nachricht ist ein JSON-\texttt{StepResult}:
\begin{lstlisting}
async with websockets.connect("ws://127.0.0.1:8001/ws") as ws:
    result = json.loads(await ws.recv())
\end{lstlisting}
\end{beispiel}

\chapter{Eingabedateiformat}
\label{app:eingabe}

Ein Netz besteht aus \textbf{fünf} DH-nahen JSON-Dateien. Der Vertrag ist
\emph{einseitig}: Die Topologie-Dateien beschreiben \emph{Trassenknoten}, nicht
das Vor-/Rücklauf-Paar, das der Solver simuliert --- der Netzbauer entfaltet das
Paar (Kapitel~\ref{ch:physik}), handgeschriebene Dateien bleiben halb so groß.
Jedes Modell verbietet unbekannte Schlüssel (\texttt{extra="forbid"}), Tippfehler
fallen also hart durch.

\begin{description}[nosep]
  \item[\texttt{network\_structure.json}] Trassenknoten (\texttt{name},
    \texttt{kind}, \texttt{geo} als \texttt{[lat, lon]}, \texttt{pn\_bar}). Der
    Bauer erzeugt aus jedem Eintrag das \texttt{\_s}/\texttt{\_r}-Junction-Paar.
  \item[\texttt{pipes.json}] Trassen (\texttt{from\_node}, \texttt{to\_node},
    \texttt{length\_km}, ISOPLUS-\texttt{std\_type} \emph{oder} explizite
    \texttt{inner\_diameter\_mm}/\texttt{u\_w\_per\_m2k}, optional
    \texttt{geometry}). Std-Typ und Einzelparameter schließen sich aus.
  \item[\texttt{consumers.json}] die Profil-Zeilen \emph{sind} die
    Übergabestationen: \texttt{q\_sh\_w}/\texttt{q\_dhw\_w} (Raum/Warmwasser)
    plus genau ein Partner: \texttt{treturn\_k} \emph{oder} \texttt{deltat\_k}
    \emph{oder} \texttt{controlled\_mdot\_kg\_per\_s}.
  \item[\texttt{producers.json}] genau ein Druck-Slack
    (\texttt{circ\_pump\_const\_pressure}, mit optionaler \texttt{heating\_curve}
    und \texttt{dp\_control}); Sekundäre als \texttt{heat\_exchanger} oder
    \texttt{pump\_mass}.
  \item[\texttt{weather.json}] \texttt{t\_amb\_c} + \texttt{t\_ground\_c}, je
    Länge \texttt{steps}.
\end{description}

Beim Laden akkumulieren alle Kreuz-Prüfungen (Knotenverweise, Array-Längen,
Horizont-Ausrichtung, genau-ein-Slack, Erreichbarkeit vom Slack,
Nullfluss/Sackgassen) in \emph{eine} Fehlermeldung. Die vollständige
Vertragsbeschreibung und die Konvertierungsentscheidungen der Referenznetze
stehen in \api{docs/REFERENCE\_NETWORKS.md}.

\begin{beispiel}[Minimalnetz: Zentrale, eine Trasse, ein Haus]
\begin{lstlisting}
// network_structure.json
{ "name": "Minimal",
  "junctions": [
    { "name": "Zentrale", "kind": "plant",    "geo": [48.0, 7.8], "pn_bar": 6.0 },
    { "name": "Haus",     "kind": "consumer", "geo": [48.0, 7.81], "pn_bar": 6.0 } ] }

// pipes.json
{ "pipes": [ { "from_node": "Zentrale", "to_node": "Haus",
               "length_km": 0.2, "std_type": "ISOPLUS_DRE50_STD" } ] }

// consumers.json  (steps Werte je Profil; hier angedeutet)
{ "resolution_minutes": 1, "steps": 1440,
  "consumers": [ { "node": "Haus", "q_sh_w": [8000, ...],
                   "q_dhw_w": [1000, ...], "treturn_k": [328.15, ...],
                   "q_design_w": 12000 } ] }

// producers.json  (der Druck-Slack)
{ "producers": [ { "kind": "slack", "node": "Zentrale",
                   "p_flow_bar": 6.0, "plift_bar": 2.0, "t_flow_k": 358.15 } ] }

// weather.json
{ "resolution_minutes": 1, "steps": 1440,
  "t_amb_c": [2.0, ...], "t_ground_c": [8.0, ...] }
\end{lstlisting}
\end{beispiel}

\chapter{Fehlerbehebung}
\label{app:fehler}
\label{sec:fehler}

\begin{description}
  \item[Port 8001 oder 5174 belegt] Der Starter erkennt einen bereits laufenden
    rtheatflow und startet ihn nicht doppelt. Eine fremde Anwendung auf dem Port
    lässt sich mit \texttt{stop\_rtheatflow.bat} räumen --- es beendet die
    Serverfenster, die Besitzer der Ports 8001/5174 und verwaiste
    node-/esbuild-/python-Prozesse dieses Repos (erkannt über die
    \emph{Kommandozeile}, nicht den Exe-Pfad). Ein parallel laufendes
    netzsim/rtpowerflow (8000/5173) bleibt unberührt --- das ist der Sinn des
    Geschwister-Portschemas.
  \item[Der Start hängt vor dem Binden] Meist der numba-Kaltstart
    (Abschnitt~\ref{sec:kaltstart}): pandapipes/numba importieren und der erste
    JIT-Lauf dauern zusammen mehrere Sekunden bis knapp eine Minute. Ein echter
    Portkonflikt ist dagegen ein \emph{schneller} \texttt{WinError 10048} mit
    sofortigem Prozessende, kein Hängen.
  \item[\api{/state} liefert 404] Normal direkt nach dem Start --- der erste
    Schritt ist noch nicht gelöst. Nach spätestens einer Schrittdauer (plus
    Kaltstart) verschwindet der Fehler von selbst.
  \item[Solver \glqq degradiert\grqq{} oder \glqq gescheitert\grqq] Die Kopfzeile
    vermerkt den Status; die Simulation heilt sich im nächsten Schritt selbst
    (Nichtkonvergenz \emph{ist} Daten, kein Absturz). \glqq degradiert\grqq{}
    heißt: Stufe~4 der Retry-Ladder (sequentiell) hat gelöst, aber die Sollwerte
    sind nicht gehalten. Häufigste Ursache für \glqq gescheitert\grqq: eine
    extreme Lastpolitik oder ein stark vermaschtes Netz --- didaktisch oft genau
    der Punkt, an dem das Netz physikalisch am Ende ist. Verbier braucht
    regulär Tier~2.
  \item[Die Karte liegt im Wasser] War ein Anker-Bug und ist behoben:
    DESTEST und Verbier tragen nur lokale Meter ohne CRS, ihre Geokoordinaten
    sind synthetische Anker. Frühere Versionen ankerten über offenem Wasser
    (Bodensee/Lac Léman); jetzt liegen die Anker an Land. Wer eigene Netze aus
    quellenlosen Metern konvertiert, ankert sie ebenfalls an Land.
  \item[Skriptzugriff auf den Vite-Dev-Server scheitert] Der Vite-Server bindet
    unter Windows nur IPv6 (\texttt{::1}). Browser sind fein (localhost löst
    zuerst nach \texttt{::1} auf), aber ein IPv4-\texttt{127.0.0.1}-curl gegen
    die Dev-UI wird abgewiesen --- für Skripte \texttt{http://[::1]:5174}
    verwenden. (Das \emph{Backend} auf 8001 ist über IPv4 erreichbar.)
  \item[Grafana zeigt nichts] Der Collector wartet, bis rtheatflow und InfluxDB
    gesund sind und der erste Schritt gelöst ist; Zeitbereich auf \glqq letzte
    5~Minuten\grqq{} stellen.
  \item[Die Schätzung erscheint nicht oder springt] Ist der Beobachter
    deaktiviert (\api{POST /estimation/config} mit \texttt{enabled=false}), fehlt
    die \texttt{estimated}-Ebene; die Sicht fällt dann auf Realität (bzw.\ im
    strikten Modus auf Gemessen) zurück. Auf großen Netzen zieht die Schätzung
    wegen der Selbstdrosselung in Sprüngen nach --- sichtbare Rechenzeit, kein
    Fehler (Abschnitt~\ref{sec:schaetzguete}).
\end{description}

\vfill
\begin{center}\small
\textit{rtheatflow --- MIT-Lizenz für Quellcode und Handbuch. Simulationskern:
pandapipes 0.14.0 (Fraunhofer IEE / Universität Kassel, BSD-3).
Schwesterprojekt für Stromnetze:} \url{https://github.com/markisbell/rtpowerflow}\textit{.}
\end{center}

\end{document}
